\documentclass[twoside,openright,titlepage,headings,footinclude=true,cleardoublepage=empty,BCOR=5mm,11pt,a4paper,english]{book}

\PassOptionsToPackage{utf8}{inputenc}
 \usepackage{inputenc}

\PassOptionsToPackage{T1}{fontenc} 
 \usepackage{fontenc}

\PassOptionsToPackage{english}{babel}
 \usepackage{babel}

\PassOptionsToPackage{autostyle}{csquotes}
 \usepackage{csquotes}

\PassOptionsToPackage{english}{varioref}
 \usepackage{varioref}

\PassOptionsToPackage{fleqn}{amsmath}
 \usepackage{amsmath}

\usepackage{siunitx,amssymb,braket,mhchem}
\sisetup{separate-uncertainty=true, multi-part-units = single}

\usepackage[dvipsnames]{xcolor}%colors più fighi
\definecolor{halfgray}{gray}{0.55}

\usepackage{geometry,appendix}

\usepackage{tabularx} % better tables
	\setlength{\extrarowheight}{3pt} % increase table row height

\usepackage{caption}
\captionsetup{format=hang,font=small,labelfont={sf,bf}}

\usepackage{subfig,multirow,booktabs}

\usepackage{bm}

\usepackage{wrapfig}

\pdfcompresslevel=9
\pdfadjustspacing=1 
\PassOptionsToPackage{pdftex}{graphicx}
	\usepackage{graphicx}
	
\usepackage{comment}

\PassOptionsToPackage{bibstyle=numeric,citestyle=numeric-comp,backend=bibtex}{biblatex}
 \usepackage{biblatex}
\addbibresource{bibliography.bib}

\usepackage{chemformula}


\begin{document}

\begin{titlepage}
\vspace{5mm}
\begin{figure}[hbtp]
\centering
\includegraphics[scale=.13]{Images/UnipdLogo.png}
\end{figure}
\vspace{5mm}
\begin{center}
{{\huge{\textsc{\bf UNIVERSIT\`A DEGLI STUDI DI PADOVA}}}\\}
\vspace{5mm}
{\Large{\bf Dipartimento di Fisica e Astronomia ``Galileo Galilei''}} \\
\vspace{5mm}
{\Large{\textsc{\bf Corso di Laurea Magistrale in Fisica}}}\\
\vspace{15mm}
{\Large{\textsc{\bf Tesi di Laurea}}}\\
\vspace{15mm}
\begin{spacing}{3}
{\LARGE \textbf{Development and characterization of an atmospheric plasma source for non-thermal blood coagulation​}}\\
\end{spacing}
\vspace{5mm}
\end{center}

\vspace{10mm}
\begin{spacing}{1.5}
\begin{tabular}{ l  c  c c c  cc c c c c  l }
{\Large{\bf Relatore}} &&&&&&&&&&& {\Large{\bf Laureando}}\\
{\Large{\bf Prof. Emilio Martines}} &&&&&&&&&&& {\Large{\bf Davide Mancini}}\\
{\Large{\bf Correlatore}}\\
{\Large{\bf Dr. Gianluca De Masi}}\\
{\Large{\bf Dr. Luigi Cordaro}}\\
\end{tabular}
\end{spacing}
\vspace{30 mm}

\begin{center}
{\Large{\bf Anno Accademico 2018/2019}}
\end{center}
\end{titlepage}
\clearpage{\pagestyle{empty}\cleardoublepage}

\thispagestyle{plain}
\begin{center}
    \vspace{2cm}
    \Large
    \textbf{Abstract}
\end{center}
In this thesis is developed and charachterized Plasma Coagulation Controller (PCC): a cold atmospheric plasma source designed for acceleration of non thermal blood coagulation. PCC is based on the Dielectric Barrier Discharge (DBD), it ionizes a neutral gas through the application of a fast voltage pulse to an electrode covered in dielectric material. Electrical charachterization of the prototype developed in this work allows to observe voltage pulses with amplitude up to \SI{10}{\kilo\volt} while current intensity is always lower than \SI{10}{\milli\ampere}, for repetition rates from \num{1} to \SI{40}{\kilo\hertz}. A peculiar phenomenon in DBD sources is that plasma propagates producing \emph{bullets}: localized zone of plasma expelled at velocities higher than \SI{10}{\kilo\meter/\second}. In this work plasma bullets are observed with a fast frame acquisition setup that allows to observe the discharge with an integration time of \SI{15}{\nano\second}. The principal analysis revolves around production and propagation of plasma bullets with different experimental setups such as: voltage pulse amplitude, composition of neutral gas (chosen between helium, neon and argon) or typology of the target placed in front of the source (e.g. insulating or conductive targets). The comparision of bullets produced with helium and neon gas suggests that bullet propagation velocity is related to electron mobility. The measurements of current intensity in a conductive target allows to observe that current flows in the target before impact with the bullet. Bullets are not observed using argon, with this gas plasma dynamics is characterized by the production of filaments that differs from bullets. 
Emission spectroscopy analysis is performed assuring the presence of \ce{OH}, \ce{N_2} and \ce{N_2^+} molecules, with different intensities for different operational regimes. The relative emission of gas components is studied for different plasma positions and gas composition, finding that during plasma propagation in air emission is always dominated by nitrogen transitions.
Temperature increase due to plasma application on an inorganic target is measured with a thermal camera. Power deposition of a single voltage pulse is found to be inversely proportional to pulse repetition rate and to target distance, while it increases with higher voltage peaks.

\clearpage


\tableofcontents

\pagestyle{headings}


\chapter{Introduzione}

\section{Plasma medicine: trattamenti al plasma su tessuti biologici}
\subsection{Coagulazione del sangue}
\section{Cold Atmospheric Plasma}
\subsection{Dielettrical Barrier Discharge}

\label{ch:intro}



\clearpage
\chapter{Electric characterization}
\label{ch:electric} 
The development of electric components used to produce the Plasma Coagulation Controller is highly influenced by the need of flexibility in settings and mobility for wound treatment. To produce plasma as DBD, in air with Helium or Argon as ignition gasses, it is necessary to produce high voltage, to permit easy medical application the design of the head must be compact with particular attention to electric safety measuraments.
The scheme outputs a voltage pulse with an amplitude up to $\SI{10}{\kilo\volt}$ and frequencies up to $\SI{50}{\kilo\hertz}$.

A rappresentation of power and signal line is in figure \ref{fig:electricline}. The circuit is divided mainly in two parts %as described in chapter \ref{ch:source}
: the controller, with alimentation and settings controls, and the head, where the discharge happens and plasma is emitted.
\begin{figure}
 \centering
 \includegraphics[width=0.8\textwidth]{Images/Linea_elettrica.png}
 \caption{Scheme of the general electric line to produce high voltage, the controller on the left and the head on the left.}
 \label{fig:electricline}
\end{figure}

The line divides in:
\begin{itemize}
 \item \bf{Alimentation} : the $\SI{220}{\volt}$ DC power line goes in a transformer that gives a $\SI{22}{\volt}$ tension to the head, passing through a diode bridge. This tension aliments the driver circuit on the head.
 \item \bf{Arduino and trigger} : the power line is reduced to $\SI{12}{\volt}$ necessaries to aliment an Arduino controller and a laser. From an analogical output of the Arduino a PWM wave goes to the laser trigger, it transmit information on the wave frequency and duration with an optical fiber that ends with a photodiode installed on the driver circuit. Wave frequency is setted by the Arduino, wave duration is setted giving the opengin time of the MOSFET installed on the driver circuit. With this setup the High-Voltage line is entirely decopuled from the controller, in this way there are not serious problems of signal reflection on the power line or the Arduino.
 \item \bf{Head} : the Driver Circuit receives a power line of $\SI{22}{\volt}$ and an optical trigger that defines frequency and duration of the voltage pulse. When the trigger gives the start signal the transformer on the head receives on primary circuit a voltage of hundreds $\si{\volt}$ and outputs from secondary circuit a voltage of thousends $\si{\volt}$. Connected to the ouput there is an electrode inside a capillary tube.
\end{itemize}


A simulation of the different signals, obteined with a simplified scheme with Spice, is in figure \ref{fig:signals}. As in the figure...

The characterization of the electric scheme is made measuring output tension and current with different settings for the pulse.


\begin{comment}
La sorgente di plasma funziona tramite l'applicazione di un'alta differenza di potenziale tra gli elettrodi separati da materiale dielettrico, come descritti nel capitolo \ref{ch:sorgente}. La tensione in ingresso nel circuito  della testa della sorgente vengono amplificati fino a tensioni di alcuni \si{\kilo\volt}. Dall'arduino di controllo è possibile regolare il tempo di apertura del circuito in un range di [$2$-$20$] \si{\micro\second} e la frequenza di lavoro in un range di [$2$-$60$] \si{\kilo\hertz}. I parametri importanti per caratterizzare il funzionamento della sorgente saranno quindi tensione e corrente all'uscita del circuito secondario del trasformatore, al variare dei parametri di funzionamento. Utilizzando una lastra metallica posta a potenziale, è possibile misurare la corrente in un dato range temporale. È inoltre possibile stimare la corrente efficace che attraversa il bersaglio, importante nel valutare gli effetti dell'applicazione del plasma.

\section{Setup delle misure di tensione e corrente}
Si vogliono effettuare misure di tensione della sorgente sia senza immissione di elio, senza formazione della plume di plasma, sia nelle condizioni di funzionamento tipiche, con flusso di elio di $\SI{2}{\liter/\minute}$.
Le tensioni vengono misurate tramite una sonda \emph{...} con attenuazione $\times1000$, le correnti tramite una sonda \emph{Tektronix CT2} che per una corrente di \SI{1}{\milli\ampere} restituisce un segnale di \SI{1}{\milli\volt}. I dati vengono letti su un oscilloscopio \emph{Yokogawa DL9040}, che permette il salvataggio dell'intera forma d'onda misurata nei diversi canali. 
Viene effettuata la caratterizzazione elettrica di entrambi i prototipi di sorgente, per i quali il circuito utilizzato è lo stesso, quindi non si aspettano variazioni significative.

\subsection{Misure senza elio}
La sonda ad alta tensione viene collegata all'uscita del circuito secondario, mentre una sonda con attenuazione $\times10$ viene utilizzata per controllare il segnale in ingresso.
La sorgente viene azionata variando la frequenza ($f$) e duty cycle ($\Delta t$). Scelta la frequenza di lavoro, viene variata la duty cycle in un range utile, considerando il tempo necessario al terminare delle oscillazioni del segnale prima dell'arrivo di una nuova onda quadra (per frequenze maggiori si potrà arrivare a duty cycle minori).
Si ottengono curve come in figura \ref{fig:tensione_es}. La risoluzione della misura viene variata in modo da avere l'errore di misura minore possibile.

\begin{figure}
\centering
\includegraphics[width=.7\textwidth]{Immagini/tensione_es_2.png}
\caption{Esempio di misura di tensione, $f = \SI{5}{\kilo\hertz}$ e $\Delta t = \SI{16}{\micro\second}$, prototipo 1.}
\label{fig:tensione_es}
\end{figure}

\subsection{Misure con elio}
Per caratterizzare il funzionamento della sorgente nelle possibili condizioni di trattamento, vengono misurate contemporaneamente, su due diversi canali dell'oscilloscopio, tensione alla quale si trova l'elettrodo e corrente che fluisce nel plasma. Per la misura di corrente viene fatto impattare il plasma su una piastra di rame  di dimensioni \SI{2}{\centi\metre} $\times$ \SI{2}{\centi\metre} $\times$ \SI{0.5}{\centi\metre}, ad una distanza di \SI{1}{\centi\metre} dall'elettrodo della sorgente, collegata alla sonda di corrente.
Nuovamente vengono effettuate misure al variare di frequenza ($f$) e duty cycle ($\Delta t$), con le stesse modalità delle misure di tensione.
Tutte le misure sono effettuate con flusso di gas He pari a $\SI{2}{\litre/\minute}$.
Si ottengono curve come in figura \ref{fig:corrente_es}.

\begin{figure}
\centering
 \includegraphics[width=.7\textwidth]{Immagini/corrente_es.png}
\caption{Esempio di misure di corrente per $f = \SI{5}{\kilo\hertz}$ e $\Delta t = \SI{24}{\micro\second}$, prototipo 1.}
\label{fig:corrente_es}
\end{figure}


\section{Presentazione misure ed analisi}
Sia per le misure di tensione che per le misure di corrente si trova un picco negativo come presentato nelle Figure \ref{fig:tensione_es} e \ref{fig:corrente_es}. Il picco di tensione ha valori tipici tra i $\num{3}$ e i $\SI{10}{\kilo\volt}$ in assenza o in presenza di gas, mentre quello di corrente tra i $\num{2}$ e i $\SI{12}{\milli\ampere}$.
L'analisi dei dati prevede la ricerca del massimo della tensione e della corrente nelle diverse configurazioni.

\subsection{Tensione di picco senza elio}
L'andamento medio delle misure presenta un picco negativo pronunciato, compatibile con i tempi di apertura del circuito. Dato un set, il valore di picco viene cercato calcolando la trasformata di Fourier del segnale (tramite le routine fftw3 delle librerie ROOT), tagliando le oscillazioni ad alta frequenza e ricostruendone una media. Nel segnale medio così ricostruito il valore del minimo viene trovato interpolando con una funzione di Landau attorno il minimo, in modo da riprodurre l'asimmetria del picco.
A queste misure viene aggiunto l'errore dovuto al taglio delle alte frequenze, preso come una media del valore assoluto dell'oscillazione del segnale tagliato. Viene inoltre aggiunto l'errore caratteristico dello strumento di misura, trascurabile rispetto l'errore dovuto alle oscillazioni veloci.
In figura \ref{fig:landau} un esempio del fit.

\begin{figure}
\centering
\includegraphics[width=.6\textwidth]{Immagini/esfiterr.png}
\caption{Esempio di fit del picco di tensione, $f = \SI{5}{\kilo\hertz}$ e $\Delta t = \SI{16}{\micro\second}$}
\label{fig:landau}
\end{figure}


Le tensioni del picco così calcolate, al variare della duty cycle per le diverse frequenze, sono presentate in figura \ref{fig:tensioni}.
Per tutte le frequenze di lavoro tra i $\SI{4}{\micro\second}$ e i $\SI{16}{\micro\second}$ risulta un andamento lineare, con tensione variabile tra i $\SI{2}{\kilo\volt}$ e i $\SI{9}{\kilo\volt}$. Aumentando ancora il tempo di apertura del circuito la tensione arriva a valori più elevati, fino un massimo di circa $\SI{10}{\kilo\volt}$, ma viene perso l'andamento lineare.


 \begin{figure}
\centering
\includegraphics[width=.6\textwidth]{Immagini/nogas.png}
\caption{Tensioni al variare del tempo di apertura del circuito e per diverse frequenze.}
\label{fig:tensioni}
\end{figure}


Le misure non sembrano presentare un andamento in funzione della frequenza, per verificarlo vengono calcolati i coefficenti dell'interpolazione lineare per le varie frequenze, presentati in figura \ref{fig:fitlin}.
Viene confermata l'assenza di un andamento specifico al variare della frequenza.

\begin{figure}
\centering
\subfloat[][Pendenza delle rette.]
  {\includegraphics[width=.48\textwidth]{Immagini/m_freq_nogas.png}}
\subfloat[][Intercetta delle rette.]
  {\includegraphics[width=.48\textwidth]{Immagini/q_freq_nogas.png}}
\caption{Parametri dell'interpolazione lineare dei set di misure senza immissione di gas(nel range $\Delta t$ stabilito), al variare della frequenza.}
\label{fig:fitlin}
\end{figure}


\subsection{Tensione e corrente con elio}
Le misure di tensione presentano l'andamento trovato precedentemente, mentre le misure di corrente presentano un primo picco negativo seguito da un picco più basso di segno opposto, positivo. L'analisi proposta è uguale a quella pensata per i set di misure precedenti: vengono tagliate le oscillazioni ad alta frequenza, ricostruito il segnale (aggiungendo l'errore dovuto al taglio delle alte frequenze e agli strumenti di misura) e il valore del minimo viene trovato interpolando con una funzione di Landau. Da questo fit vengono calcolati valori e posizione del picco di tensione, del picco negativo di corrente e del picco positivo di corrente.
In figura \ref{fig:landau} un esempio del fit.

\begin{figure}
\centering
\includegraphics[width=.6\textwidth]{Immagini/es_fit_landau.png}
\caption{Esempio di fit del picco di corrente, $f = \SI{5}{\kilo\hertz}$ e $\Delta t = \SI{24}{\micro\second}$}
\label{fig:landau}
\end{figure}

Le tensioni del picco così calcolate, al variare della duty cycle per le diverse frequenze, sono presentate in figura \ref{fig:picchi}.
Nuovamente troviamo un andamento lineare per la tensione tra i $\num{4}$ e i $\SI{16}{\micro\second}$. Anche il picco negativo di corrente presenta questo andamento lineare, mentre per il picco positivo non è possibile identificare un comportamento simile, i valori si disperdono.

\begin{figure}
\centering
\subfloat[][Modulo del picco di tensione.]
  {\includegraphics[width=.65\textwidth]{Immagini/vpp_corrente.png}}
\\
\subfloat[][Modulo del picco primario di corrente.]
  {\includegraphics[width=.65\textwidth]{Immagini/I1_corrente.png}}
\\
\subfloat[][Picco secondario di corrente.]
  {\includegraphics[width=.65\textwidth]{Immagini/I2_corrente.png}}
\caption{Valori dei picchi misurati al variare del tempo di apertura del circuito e per diverse frequenze.}
\label{fig:picchi}
\end{figure}


Nuovamente, per visualizzare in maniera esplicita l'effetto della variazione della frequenza, vengono calcolati i coefficenti dell'interpolazione lineare per le tensioni di picco e per le correnti di picco negativo, presentati in figura \ref{fig:fitlin_cor}. 
Per le tensioni risulta un comportamento identico al precedente, dove i valori si assestano attorno una media lievemente inferiore rispetto le misure in assenza di elio.
Per il valore massimo di corrente viene trovato un aumento in funzione della frequenza di funzionamento della sorgente.

\begin{figure}
\centering
\subfloat[][Pendenze dei picchi di tensione.]
  {\includegraphics[width=.48\textwidth]{Immagini/mVpp_cor.png}}
\subfloat[][Intercette dei picchi di tensione.]
  {\includegraphics[width=.48\textwidth]{Immagini/qVpp_cor.png}}
\newline
\subfloat[][Pendenze dei picchi di corrente.]
  {\includegraphics[width=.48\textwidth]{Immagini/mI1_cor.png}}
\subfloat[][Intercette dei picchi di corrente.]
  {\includegraphics[width=.48\textwidth]{Immagini/qI1_cor.png}}
\caption{Parametri dell'interpolazione lineare dei picchi delle misure di tensione e corrente.}
\label{fig:fitlin_cor}
\end{figure}

\begin{comment}

Data la possibilità di visualizzare contemporaneamente sia la tensione sia la corrente in uscita dal circuito, viene proposta un'analisi delle variazioni temporali tra i diversi picchi. In particolare viene calcolato il tempo tra i due picchi di corrente e tra il picco di tensione e il picco di corrente primario, mostrati in figura \ref{fig:tempi}.
In entrambi i casi non è possibile estrapolare un andamento particolare, indicando che non vi sono differenze significative nei tempi di salita dei picchi al variare del tempo di apertura del circuito e della frequenza.

\begin{figure}
\centering
\subfloat[][Tempo tra i picchi di corrente.]
  {\includegraphics[width=.48\textwidth]{Immagini/ti2ti1.png}}
\subfloat[][Tempo tra picco tensione e primo picco di corrente.]
  {\includegraphics[width=.48\textwidth]{Immagini/ti1tv.png}}
\caption{Misura delle differenze temporali tra i picchi.}
\label{fig:tempi}
\end{figure}

\subsection{Corrente efficace}
Durante l'applicazione del plasma su tessuti vivi bisogna considerare il tempo di reazione effettivo del bersaglio, l'impulso di corrente alle frequenze di lavoro della sorgente presenta un periodo inferiore rispetto questi tempi. Per valutare gli effetti del trattamento viene calcolato il valore della corrente efficace che fluisce sulla piastra bersaglio in un tempo di \SI{1}{\milli\second}, dell'ordine di grandezza dei tempi di risposta da considerare (vedi articolo?), utilizzando la formula in \ref{eq:ieff}.
\begin{equation}
 \centering
 I_{\text{eff}} = \frac{1}{(t_2-t_1)}\sqrt{\int_{t_1}^{t_2} I^2 \,dt}
 \label{eq:ieff}
\end{equation}

In Figura \ref{fig:correff} vengono presentati i valori della corrente efficace in maniera simile a quanto fatto per le misure di corrente precedentemente.
A parità di tempo di apertura del circuito, una maggiore frequenza implica che nel tempo scelto di \SI{1}{\milli\second} vi sarà un numero di periodi maggiore, aumentando la corrente efficace nel circuito. In figura si vede come mediamente la corrente efficace sia più grande a frequenze maggiori, ma assume sempre valori inferiori ai \SI{3}{\milli\ampere}.

\begin{figure}
\centering
\includegraphics[width=.7\textwidth]{Immagini/ieff.png}
\caption{Corrente efficace calcolata al variare del tempo di apertura del circuito e per diverse frequenze.}
\label{fig:correff}
\end{figure}

\end{comment}

\clearpage
\chapter{Shape and dynamics}
\label{ch:forma}
When it's seen by naked eye, plasma sources for medical uses expel a little plume of plasma that emits radiation in a visible range, with different colours or intensity that depends on the gas composition, it's flow and intensity of applied electric field. 
However if we see this plasma at a specific time, recent studies shows that what is expelled is not a continuos flow, but the ``plume'' is formed by compact bullets with high velocities (\cite{Mericam_Bourdet_2009}, \cite{doi:10.1002/ppap.200900078}).
An example of this phenomenon is shown in figure \ref{fig:pl_bullet}, measured with the experimental apparatus that will be explained later.
\begin{figure}
 \centering
 \includegraphics[width=0.8\textwidth]{Images/Shape/frames.png}
 \caption{Example of plasma bullet expulsion, as measured with our experimental setup, with an helium flow of \SI{2}{\liter/\minute}.}
 \label{fig:pl_bullet}
\end{figure}

Plasma bullets still needs to be studied in depth, we only know the basic dynamic of formation and expulsion. Some general features are that:
\begin{itemize}
 \item bullet's velocities are $> \SI{10}{\kilo\meter/\second}$;
 \item bullet formation, it's velocity and it's travel distance depend on applied tension on the electrode;
 \item frontal image of the bullet, on a plane perpendicular to it's velocity, is not a full circle but it's donut shaped, with an outer ring of higher intensity.
\end{itemize}

The scope of this experiment is to observe plasma bullets produced with our source, their shape and their velocity and how they change with different discharge conditions.

Given their typical velocities and the temperature of the plasma, bullet propagation is tought to be related to a travelling ionization front. This propagation can be studied with simplified simulation of DBD discharges, where it's reproduced the behaviour of plasma bullets (\cite{doi:10.1063/1.4963115}, \cite{Breden_2012}) or the interaction between plasma and a target (\cite{doi:10.1063/1.4923345}).
%We present a simple 1D model that can reproduce the propagation of this ionization front.

\section{Experimental setup}
To visually observe dynamics of plasma formation and propagation, it is needed an acquisition setup with an high-speed camera that has little integration time, around \SI{10}{\nano\second}, and an image intensifier that permits to visualize light emitted in such a short time interval.

This is a measure where we need coincidence between discharge and frame measure so it is necessary to consider instruments and plasma source specific delays and give appropriate triggers.

Experimental setup is shown in photo \ref{fig:fotosetup} and a scheme is presented in \ref{fig:schemashape}, there are the trigger signal lines, optical acquisition pointed at source exit and measure acquisition composed by a computer and an oscilloscope. 
\begin{figure}
 \centering
 \includegraphics[width=0.6\textwidth]{Images/Shape/acq_ottica.png}
 \caption{Photo of the experimental setup.}
 \label{fig:fotosetup}
\end{figure}

\begin{figure}
 \centering
 \includegraphics[width=0.7\textwidth]{Images/Shape/acq_ottica.png}
 \caption{Experimental setup scheme. Function generator 1, function generator 2 and delay generator send trigger signal to camera (CCD), to source and to image intensifier (MCP), full lines in the scheme. Camera sends measured frame to a computer (PC) and from source, source's target, function generator 2 and delay are taken signals read on the oscilloscope, pointed lines in the scheme.}
 \label{fig:schemashape}
\end{figure}



\subsection{Source and optical setup}
Source head is positioned vertically and emitted light is collected from the side.
Optical apparatus is composed by a lens %specific
 coupled with a Micro Channel Plate image intensifier (MCP). An MCP works as in figure \ref{fig:MCP}: for every photon received, thanks to an high-voltage power supply, it emits many photons with little angular deviation.
\begin{figure}
 \centering
 \includegraphics[width=0.4\textwidth]{Images/Shape/MCPsingle-stage.jpg}
 \caption{Micro Channel Plate image intensifier functioning.}
 \label{fig:MCP}
\end{figure}

Those photons are received by an high-speed camera \emph{Point Grey Flea} (\cite{flea}) equipped with a CCD of $\num{1024} \times \num{768}$ square pixels $\SI{4.75}{\micro\meter}$ wide. Every frame is sent to the pc where FLIR software elaborates and saves them in pgm format (\cite{pgm}).

\subsection{Source, power lines and electric signals}
It's used a source with electric features and functioning similar to those described before, with an output between $\num{3}$ and \SI{8}{\kilo\volt}, variabile changing $\Delta t$ of the trigger signal (see \ref{ch:electric}), that is always given through an optical fiber with frequency $f$. Differances are that tension is positive and not negative, because it helps formation of the plume with gasses harder to ignite (e.g. Argon), and that trigger signal is given by a generator function ($2$ in the scheme) to define discharge time respect to the trigger of other instruments.

The source is positioned vertically at a distance of \SI{42.0(1)}{\milli\meter} from optical bench, with the glass nozzle that permits to observe plasma formation inside it (external diameter \SI{8.0(1)}{\milli\meter}, internal diameter \SI{6.0(1)}{\milli\meter}), with a distance from the end of the electrode to nozzle exit of \SI{12}{\milli\meter}. At the exit the nozzle shrinks for \SI{3.0(1)}{\milli\meter}, until a diameter of \SI{5.0(1)}{\milli\meter}

Under the source is possible to put targets. Are used two different targets at different heights: conductive target and insulator target. The first one is a copper square sheet of dimensions $\SI{10}{\milli\meter} \times \SI{10}{\milli\meter} \times \SI{1}{\milli\meter}$ (used for current measures in chapter \ref{ch:electric}), the second one is a simple plastic material.


CCD camera is powered by a simple alimentator while image intensifier MCP is powered by an high-voltage supply, both are triggered by the delay generator, with different times (see next section).


On an oscilloscope are measured the trigger signal given to source head, the trigger signal given to MPC, the voltage electrode with high-voltage probe \emph{Tektronix P6015A} and the current intensity when it's used the conductive target. Current intensity measure is done measuring voltage drop on a resistance of $\SI{1}{\mega\ohm}$ with a probe $x\num{10}$.

\subsection{Trigger synchronization}
Experiment's scope is to know plasma dynamics at a specific point on voltage waveform, so it's necessary to know precisely discharge and measure times.

To compose the necessary trigger line are utilized:
\begin{itemize}
 \item function generator \emph{Or-x 310}, $1$ in figure, that sends a square wave with set amplitude, width and frequency $f$;
 \item function generator \emph{Lecroy 9210}, $2$ in figure, that sends a square wave with frequency given by the trigger, amplitude set, and variable width $\Delta t$;
 \item delay time generator \emph{Stanford DG535}, that sends square wave with set amplitude, frequency given by the voltage input and starting times given by voltage input startin time plus settable delays (4 channels).
\end{itemize}

Every intrument has it's own time that elapse between trigger signal and effective measure. The longer one is arming time for high-speed camera in the order of \si{\milli\second}, the smaller one is integration time for acquisition system, that starts from the activation of the image intensifier and span $\SI{15}{\nano\second}$.

A time line is shown in figure \ref{fig:times}, an example of signal taken with the oscilloscope is in figure \ref{fig:times_signals}. There are three relevant times defined by function and delay generators:
\begin{enumerate}
 \item $t_{0}$ is the starting time for the square wave given by \emph{Function generator 1}, with an amplitude of $\SI{5}{\volt}$ and a frequency $f$, that goes as external trigger to \emph{Function generator 2} and as voltage input to \emph{Delay generator}. From \emph{Function generator 2}, at $t_{0}$, starts the square wave that triggers discharge, with a width of $\Delta t$ that will define voltage amplitude (see chapter \ref{ch:electric}) and the same frequency $f$. From \emph{Delay generator}, at $t_{0}$, without delays, starts the trigger signal to arm the camera with the CCD.
 \item $t_{\text{DIS}}$ is the effective discharge time, when the amplitude peak starts. From trigger signal end to amplitude peak start there is a time delay given mainly by the response time of photodiode. Measuring the signals as in figure \ref{fig:times_signals}, it's possible to estimate this delay as $\SI{987.7(567)}{\nano\second}$, constant for every $f$ and $\Delta t$. Once the discharge starts, in the grey zone in figure, there are the events that we want to measure, plasma formation and propagation.
 \item $t_{\text{MIS}}$ is the measure time, when the MCP is triggered on. The delay between $t_{0}$ and $t_{\text{MIS}}$, $t_{D}$, is given by the \emph{Delay generator} with possible steps of \SI{1}{\pico\second}. Changing $t_{D}$ during measure it's possible to see plasma dynamics at different times that corresponds to different points on electrode tension waveform, as in figure \ref{fig:times_signals}.
\end{enumerate}
\begin{figure}
 \centering
 \includegraphics[width=0.6\textwidth]{Images/Shape/times.png}
 \caption{Time signal synchronization scheme: $t_{0}$ is the starting trigger time, $\Delta t$ is the opening time for plasma source (see chapter \ref{ch:electric}), $t_{\text{DIS}}$ is the starting time for the discharge and $t_{\text{MIS}}$ the starting time for the MCP i.e. the measure time}.
 \label{fig:times}
\end{figure}

\begin{figure}
 \centering
 \includegraphics[width=0.6\textwidth]{Images/Shape/times_osc.png}
 \caption{Oscilloscope measure example, in green the discharge trigger, in red the electrode tension output and in blue the MCP trigger.}
 \label{fig:times_signals}
\end{figure}


Integration time for a single frame is \SI{15}{\nano\second}, so the time step between two measures has to be larger then it. For slow plasma bullets (Neon or Argon) it's used a step of \SI{50}{\nano\second}, for fast bullets (Helium) it's used a step of \SI{20}{\nano\second}.


The effective measure time have to consider slower time of the camera, larger then \SI{50}{\milli\second}. It's important to note that also with little frequencies, around \SI{50}{\hertz} (see later the minimum working frequency for different gasses), time between two pulses is \SI{20}{\milli\second}, so it is not possible to see two consecutive frames, we observe the same time in different discharges.

Changing the settings on camera acquisition software is possible to measure a signal below pixel's saturation point. One of the parameters is the \emph{Shutter time}, the opening time of the camera shutter, that if set on a value larger then the time between two signals permits to integrate between two frames.
When we work with a gas, we select an appropriate \emph{Gain} value and an appropriate \emph{Shutter time}, chosing to work measuring a single discharge frame or a multiple discharge frame.

\subsection{Different setups}
Plasma formation is influenced by many parameters such as frequency of high-voltage pulses, their rise time and maximum, gas type, gas flow, presence of a target and its features (see \cite{Mericam_Bourdet_2009}, \cite{Jarrige_2010}).

For pulse frequency we find a lower limit value, different for each ignition gas, under which there is no discharge. For frequencies higher then this value we don't expect changes (as there aren't in the electric behaviour in chapter \ref{ch:electric}), so we use a single chosen value that permits plasma ignition and doesn't stress the experimental setup.
High-voltage values and gas type determine plasma bullet formation and affect its expulsion velocity. In particular, different gasses have different atomic mass and ionization energies, they will have different voltage ignition values and reactive species formed in plasma will have different velocities.
Gas flow influences how much gas there is when the ignition starts, so varying it we observe different bullets diameter and velocity.
Target influences bullet expulsion, propagation and plasma rebound signal. With a conductive target near plasma exit we will have an easier path seen by charges that propagates and, once the bullet hits the target, its observed higher luminosity going from the target to the electrode. With an insulator target we don't find the ionization channel closing on the target, but we observe charge deposition on the insulator with a shape that depends from the gas and the target.
In some setups, to help plasma ignition it's also used a conductor ring at ground potential or at floating potential, placed around source nozzle, after the electrode position.


In this study we present different setups with three different gasses, as shown in table \ref{tab:setups}.

\begin{table}
 \centering
 \begin{tabular}{cccccccc}
  \toprule
  Gas   &Setup  &$\Delta t$ &Target &Target position    &Flow rate  &Other  &$\Delta t_D$\\
  \midrule
  \multirow{8}*{\ce{He}}    &A  &$\num{30}, \num{35}, \num{40}$ &Conductor  &\SI{24}{\milli\meter}  &\SI{2}{\liter/\minute} &-  &\SI{50}{\nano\second}\\
                            &B  &$\num{30}, \num{35}, \num{40}$ &Conductor  &\SI{32}{\milli\meter}  &\SI{2}{\liter/\minute} &-  &\SI{20}{\nano\second}\\
                            &C  &$\num{30}, \num{35}, \num{40}$ &Insulator  &\SI{24}{\milli\meter}  &\SI{2}{\liter/\minute} &-  &\SI{20}{\nano\second}\\
                            &D  &$\num{30}, \num{35}, \num{40}$ &-  &-  &\SI{2}{\liter/\minute} &-  &\SI{20}{\nano\second}\\
                            &E  &$\num{35}$ &-  &-  &\SI{2}{\liter/\minute} &Ground ring  &\SI{20}{\nano\second}\\
                            &F  &$\num{35}$ &-  &-  &\SI{1}{\liter/\minute} &-  &\SI{20}{\nano\second}\\
                            &G  &$\num{35}$ &-  &-  &\SI{3}{\liter/\minute} &-  &\SI{20}{\nano\second}\\
                            &H  &$\num{35}$ &-  &-  &\SI{4}{\liter/\minute} &-  &\SI{20}{\nano\second}\\
  \midrule
  \multirow{5}*{\ce{Ne}}    &A  &$\num{20}, \num{25}, \num{30}$ &Conductor  &\SI{24}{\milli\meter}  &\SI{2}{\liter/\minute} &-  &\SI{50}{\nano\second}\\
                            &B  &$\num{20}, \num{25}, \num{30}$ &Conductor  &\SI{32}{\milli\meter}  &\SI{2}{\liter/\minute} &-  &\SI{50}{\nano\second}\\
                            &C  &$\num{20}, \num{25}, \num{30}$ &Insulator  &\SI{24}{\milli\meter}  &\SI{2}{\liter/\minute} &-  &\SI{50}{\nano\second}\\
                            &D  &$\num{20}, \num{25}, \num{30}$ &-  &-  &\SI{2}{\liter/\minute} &-  &\SI{50}{\nano\second}\\
                            &E  &$\num{30}$ &-  &-  &\SI{2}{\liter/\minute} &Ground ring  &\SI{50}{\nano\second}\\
  \midrule
  \multirow{6}*{\ce{Ar}}    &A  &$\num{35}$ &-  &-  &\SI{2}{\liter/\minute} &Ground ring  &\SI{50}{\nano\second}\\
                            &B  &$\num{35}$ &-  &\SI{32}{\milli\meter}  &\SI{2}{\liter/\minute} &-  &\SI{20}{\nano\second}\\
                            &C  &$\num{35}$ &-  &-  &\SI{2}{\liter/\minute} &Floating ring  &\SI{50}{\nano\second}\\
                            &D  &$\num{35}$ &Conductor  &\SI{20}{\milli\meter}  &\SI{2}{\liter/\minute} &Ground ring  &\SI{50}{\nano\second}\\
                            &E  &$\num{35}$ &Conductor  &\SI{20}{\milli\meter}  &\SI{2}{\liter/\minute} &-  &\SI{50}{\nano\second}\\
                            &F  &$\num{35}$ &-  &-  &\SI{2}{\liter/\minute} &-  &\SI{50}{\nano\second}\\
                            &G  &$\num{35}$ &Conductor  &\SI{32}{\milli\meter}  &\SI{2}{\liter/\minute} &-  &\SI{50}{\nano\second}\\
  \bottomrule
 \end{tabular}
 \caption{Description of measure setups. In first column there is the gas; second column is the setup name; third column is voltage pulse time width, expressed as percentage of a square trigger \SI{10}{\micro\second} long; four and fifth columns are target information, where it's used a target; sixth column is used gas flow; seventh column are other informations, i.e. if it's used the conductor ring explained before; eight column is the time step between two different frames.}
 \label{tab:setups}
\end{table}


\subsection{Frame analysis and calibration}
Once a measure setup is chosen, $t_D$ is set around the start of the discharge, we measure the oscilloscope waveform for every channel (setup as explained before), and from them we save $5$ frames for every $t_D$ varying it with steps of $\Delta t_D$. Measures are taken until the end of the main tension peaks and the camera doesn't see anything.


From oscilloscope waveforms we can extrapolate the measure time from the start of the discharge and tension value at measure time. When there is the conductive target it's also possible to see current intensity and it's peak time.


Frame analysis is done converting the pgm files in 2-dimensional histograms, using \emph{TH2} class written in \emph{ROOT} libraries (see \cite{ROOT:TH2}).
We are interested in plasma bullet formation and it's expulsion, but also in other fenomenological behaviour that is observed during measures.

For plasma bullet we take the mean for every pixel between the five frames taken, from then it's isolated the plasma bullet as the collection of pixels with maximum intensity in the frame. From this image we extrapolate the position of the bullet with the coordinates of it's luminosity barycenter on the plane seen by camera, called x-y in the analysis.
Once established the center, we can find it's contour as the points were luminosity goes under a certain percentage of the maximum, from that we estimate it's dimensions along x and y axis and it's average luminosity.
With barycenter coordinates it's possible to compute their velocity using a finite differance formula (that takes into consideration possible different time steps, see \cite{Bhadauria}).


Actual dimensions on a frame are found with a calibration with a known target: it is used a plate with $4$ holes, diameter of \SI{1.0(1)}{\milli\meter} at the vertice of a square with an edge lenght of \SI{10.0(1)}{\milli\meter}, positioned at plasma exit (source is off) and illuminated from behind with a torch.
Acquiring a frame of this setup it's possible to extrapolate pixel's distance that corresponds to \SI{10}{\milli\meter} for every square edge, average them, and calculate pixel's width in our frames, resulting a value of $d_{pix} = \SI{0.172(2)}{\milli\meter}$.

\section{Helium flow}
We start the presentation of measures with helium, as it is commonly used for medical applications of plasma.
Helium it's an element with standard atomic weight of $\num{4.002}$, 1st ionization energy of \SI{24.587}{\electronvolt} and 2nd ionization energy of \SI{54.418}{\electronvolt}.
It's easy to produce helium plasma, when it's ignited it emits radiation with principal wavelenghts in violet and orange ($\lambda = \num{388.86}$ and \SI{587.56}{\nano\meter}), resulting in a purple plume at naked eye.
Principal downside of helium gasses it's that they are expensive.


In our experiment we produce helium plasma with pulse frequency $f = \SI{5}{\kilo\hertz}$ and $\Delta t = \num{3}$, $\num{3.5}$ and \SI{4}{\micro\second}, as explained in table \ref{tab:setups}.
Delays are setted to take a single frame during every acquisition, for every measure we take $5$ frames and average them.


First we present a tipical discharge without target, after is shown how the behaviour changes varying parameters and then it's discussed how a target influences bullet propagation.

\subsection{Plasma bullet description}
We use as example bullet the setup D with $\Delta t = \SI{3.5}{\micro\second}$, i.e. absence of a target and gas flow of \SI{2}{\liter/\minute}.


It's possible to divide the phenomenon that we want to observe in four phases, as presented in figure \ref{fig:bullet_es}: bullet formation, bullet displacement inside the nozzle, bullet expulsion and bullet propagation outside the nozzle, in air. It's important to remember that time values goes from the start of the voltage peak, extrapolated from oscilloscope measures.
When the voltage goes over a definite value, we observe a rapid increase in luminosity around the electrode that it's interpreted as plasma formation. From there, in a time of $\sim \SI{50}{\nano\second}$ the plasma modifies it'shape and becames what we call the bullet: a zone with high luminosity confined in space. The bullet then moves in the exit direction and it's expelled always with a well defined front. In air we see the round shaped bullet propagation, until it reaches a maximum travel distance and then luminosity decreases rapidly, until we see nothing more. In these measure conditions we see again plasma formation in correspondence of the negative tension peak, but it is  a rapid process without propagation.
\begin{figure}
 \centering
 \subfloat[\SI{140.93}{\nano\second}]{
    \includegraphics[width=0.48\textwidth]{Images/Shape/elio_d035_es1.png}
 }
 \hfill
 \subfloat[\SI{240.93}{\nano\second}]{
    \includegraphics[width=0.48\textwidth]{Images/Shape/elio_d035_es2.png}
 }
 
 \subfloat[\SI{400.93}{\nano\second}]{
    \includegraphics[width=0.48\textwidth]{Images/Shape/elio_d035_es3.png}
 }
 \hfill
 \subfloat[\SI{460.93}{\nano\second}]{
    \includegraphics[width=0.48\textwidth]{Images/Shape/elio_d035_es4.png}
 }
 \caption{Four phases of plasma dynamics: bullet formation (a), bullet displacement inside the nozzle (b), bullet expulsion (c) and bullet propagation outside the nozzle (d). Time intervals are calculated from the start of the voltage peak.}
 \label{fig:bullet_es}
\end{figure}


\paragraph{Electrode voltage and bullet intensity}
The time interval in wich we see the bullet can be defined as where we can see a definite zone with mean luminosity higher then the background. The starting time is given by voltage reaching a certain value, when we see bullet formation, and the bullet expires when it's luminosity decreases.

In figure \ref{fig:elio_d035_I} are presented the mean intensity of the plume during the motion and the zoom on the intersted time interval for voltage measure. Plasma formation is observed around \SI{120.93}{\nano\second} after peak start, at a tension value of \SI{5710.00(28550)}{\volt}. %while maximum peak value is \SI{6588.80(32944)}{\volt} at time \SI{260.93}{\nano\second}.
As we can see, luminosity is higher during plasma formation around the electrode and it decrease very little during all the propagation, even after the expulsion in air (pointed line on the graph). After a time of \SI{440}{\nano\second} the bullet luminosity decreases rapidly and it disperses in the air.
\begin{figure}
 \centering
 \subfloat[Zoom on voltage waveform.]{
    \includegraphics[width=0.48\textwidth]{Images/Shape/elio_d035_V.png}
 }
 \hfill
 \subfloat[Mean intensity of plasma bullet.]{
    \includegraphics[width=0.48\textwidth]{Images/Shape/elio_d035_Im.png}
 }
 \caption{Voltage waveform and intensity of the bullet versus time. Time zero is the start of the voltage peak, ending time is taken as the time where intensity reaches background value. The pointed line in intensity graph indicates the end of the nozzle.}
 \label{fig:elio_d035_I}
\end{figure}

\paragraph{Barycenter coordinates and direction}
The motion of the bullet on observing plane can be extrapolated from it's barycenter coordinates.

In figure \ref{fig:elio_d035_bary} there are the two coordinates versus time.
Coordinate x, along the axis of the source (of the electrode and the nozzle) is more relevant to observe bullet expulsion. We see as plasma forms around position $\SI{73}{\milli\meter}$ and the nozzle ends at $\SI{85}{\milli\meter}$, compatible with the distance electrode-nozzle exit described before. After plasma formation we can see bullet displacement in the nozzle with a certain velocity, until it reaches the exit and propagates in the air more rapidly. The bullet travels until it reaches a maximum x, covering a distance of \SI{26.08(2)}{\milli\meter} from the electrode. It seems that the bullet stops while it's luminosity decreases.

Coordinate y it's relevant to show bullet direction, as we can see, it moves in a space interval of \SI{2}{\milli\meter}. The bullet forms over the electrode, then, when it enlarges to cover all the nozzle diameter, lowers it's barycenter with a constant slope, even after the expulsion. Once the bullet starts to decrease in luminosity, it stops its motion also on the y direction.
\begin{figure}
 \centering
 \subfloat[]{
    \includegraphics[width=0.48\textwidth]{Images/Shape/elio_d035_xb.png}
 }
 \hfill
 \subfloat[]{
    \includegraphics[width=0.48\textwidth]{Images/Shape/elio_d035_yb.png}
 }
 \caption{Barycenter coordinates of the bullet. Pointed lines indicate the end of the nozzle.}
 \label{fig:elio_d035_I}
\end{figure}

Lowering of the baricenter in the y direction it's explainable by a tilt in the source, that is not perfectly perpendicular to the optical bench, if we consider the entire dimension of the bullet in the y direction, with it's upper and lower limits. In figure \ref{fig:elio_d035_diry} are shown y maximum, minimum and baricenter of the bullet, in function of the x brycenter coordinate and can be seen how the entire figure points downward (to the side turning the image). Inside the nozzle we can take y maximum as a value that defines the contour of the nozzle, with a fit as in figure we find a value of \SI{3.39(10)}{\degree}. Outside the nozzle we find that the barycenter direction is pointed \SI{3.71(4)}{\degree} downward, a little higher then the direction of the nozzle. The values are almost compatible with each other, so it seems to suggest that tilting the source it's possible to direct the bullet even after it's propagation in air.
\begin{figure}
 \centering
 \includegraphics[width=0.6\textwidth]{Images/Shape/elio_d035_incl.png}
\end{figure}


\paragraph{Bullet dimensions}
Once we define the contour of the bullet, it's simple to take it's dimensions along the two directions.

In figure \ref{fig:elio_d035_dim} are presented width and height of the bullet.
In the x direction, the bullet presents constant diameter until it approaches nozzle exit, where it enlarges a little to reach the exit. In the air the bullet mantaines its dimension, with values a little lower then inside the nozzle, but compatible within the error. When it stops it seems to have greater diameter, but the measure loses significance as the bullet fades in the air and it can be confused with the background.
In y direction we find a constant value of \SI{4.65(20)}{\milli\meter} during propagation in the nozzle, as expected because the bullet covers all the nozzle area, a little lower then its diameter due to glass refraction effects. Once the nozzle shrinks, the bullet diminish it's diameter and maintains a constant value of \SI{2.07(20)}{\milli\meter} during all the propagation in air, until it stops.
\begin{figure}
 \centering
 \subfloat[]{
    \includegraphics[width=0.48\textwidth]{Images/Shape/elio_d035_dx.png}
 }
 \hfill
 \subfloat[]{
    \includegraphics[width=0.48\textwidth]{Images/Shape/elio_d035_dy.png}
 }
 \caption{Dimensions of the bullet. Pointed lines indicate the end of the nozzle.}
 \label{fig:elio_d035_I}
\end{figure}


\paragraph{Bullet velocity}
From barycenter graphs we see as bullet motion has a definite velocity during described phases.
Fitting the graph with a line during the displacement phase inside the nozzle, as in figure \ref{fig:elio_d035_vx}, we find that it moves with a speed of $v_{N} = \SI{46.48(20)}{\kilo\meter/\second}$. Once it exits the nozzle it speeds up, until a value of $v_{A} = \SI{95.16(60)}{\kilo\meter/\second}$.


With a 3 point finite differance formula we can estimate it's velocity point by point (excluding the first and the last), finding the values shown in figure \ref{fig:elio_d035_vx}. From those values we can average velocities during the two motion phases, finding compatible mean value respect the linear fit: $v_{N} = \SI{48.23(103)}{\kilo\meter/\second}$ and $v_{A} = \SI{95.16(182)}{\kilo\meter/\second}$.
\begin{figure}
 \centering
 \subfloat[Fit of x barycenter.]{
    \includegraphics[width=0.48\textwidth]{Images/Shape/elio_d035_xb_fitv.png}
 }
 \hfill
 \subfloat[Velocities in x direction.]{
    \includegraphics[width=0.48\textwidth]{Images/Shape/elio_d035_vx_fitv.png}
 }
 \caption{Velocity of the bullet in x direction. Pointed lines indicate the end of the nozzle, red lines calculated values during displacement and propagation phases.}
 \label{fig:elio_d035_vx}
\end{figure}

It's interesting to point out that those velocities are much higher then the average velocity of the neutral gas, that for a flow of \SI{2}{\liter/\minute} is of \SI{12}{\meter\second}.

\subsection{Voltage influence}
With an higher $\Delta t$ we can reach higher voltage peak values (see chapter \ref{ch:electric}), so we will have higher electron energies. With more energy we expect that the bullet will reach furter distances, with higher velocity. In figure \ref{fig:elio_d} and table \ref{tab:elio_d} are presented the results of the analysis.

The discharge starts always in the same voltage interval, around \SI{5.7}{\kilo\volt}. We observe a higher luminosity with higher tensions, but it decreases always in a time interval around \SI{100}{\nano\second} once it exits the nozzle (figure \ref{fig:elio_d} (a)).

The plot for the barycenter coordinates (figure \ref{fig:elio_d} (b) - (c)) shows how with higher voltage the bullet reaches further distances ($x_{\text{dist}}$ and $y_{\text{dist}}$ in table). For the lower voltage peak it seems that the bullet doesn't propagate in air, as it stops right after nozzle exit.

Dimensions of the bullet (figure \ref{fig:elio_d} (d) - (e)) are not influenced by voltage value, as we can see they have the exact same behaviour, around the same values, for every opening time. A notable differance it's in the y diameter inside the nozzle, where for higher tension we have a slightly larger value. As mentioned before here we have refraction effects due to the glass nozzle, different for those measures due to the higher luminosity.

Velocity values (figure \ref{fig:elio_d} (f)) show how the bullet behaviour is solid during the propagation inside the nozzle and becames more erratic once it exits. As mentioned before, for the lowest opening time we have lower velocities and a rapid drop after contact with air, in this configuration it's not possible to find a propagation velocity. For higher opening times we find the same behaviour, with a proportionality between peak voltage value and velocities: if before we found a stable motion with $v_{A} = \SI{92.93(6)}{\kilo\meter/\second}$, with higher tension it reaches $v_{A} = \SI{149.47(9)}{\kilo\meter/\second}$.

\begin{figure}
 \centering
 \subfloat[Mean luminosity.]{
    \includegraphics[width=0.48\textwidth]{Images/Shape/elio_d_Im.png}
 }
 \hfill
 \subfloat[x barycenter.]{
    \includegraphics[width=0.48\textwidth]{Images/Shape/elio_d_xb.png}
 }
 \subfloat[y barycenter.]{
    \includegraphics[width=0.48\textwidth]{Images/Shape/elio_d_yb.png}
 }
 \hfill
 \subfloat[x diameter.]{
    \includegraphics[width=0.48\textwidth]{Images/Shape/elio_d_dx.png}
 }
 \subfloat[y diameter.]{
    \includegraphics[width=0.48\textwidth]{Images/Shape/elio_d_dy.png}
 }
 \hfill
 \subfloat[x velocity.]{
    \includegraphics[width=0.48\textwidth]{Images/Shape/elio_d_vx.png}
 }
 \caption{Results of the analysis for setup d (helium flow of \SI{2}{\liter/\minute} without target), with different peak voltage values.}
 \label{fig:elio_d}
\end{figure}

\begin{comment}
\begin{table}
 \centering
 \begin{tabular}{ccccccc}
 \toprule
 $V_{p}$ [kV]    &$t_{0}$ [ns] &$V_{0}$ [kV]    &$x_{\text{dist}}$ [mm]   &$y_{\text{dist}}$ [mm]   &$v_{N}$ [km/s]   &$v_{A}$ [km/s]\\
 \midrule
 \num{5.66(30)}  &\num{131.5}  &\num{5.55(28)}    &\num{15.99(1)} &\num{1.66(2)}  &\num{37.72(4)} &-\\
 \num{6.59(33)}  &\num{120.9}  &\num{5.71(29)}    &\num{26.08(2)} &\num{1.92(4)}  &\num{46.48(2)} &\num{95.16(6)}\\
 \num{7.22(36)}  &\num{119.5}  &\num{6.01(30)}    &\num{34.55(5)} &\num{2.07(1)}  &\num{59.80(3)} &\num{149.47(9)}\\
 \bottomrule
 \end{tabular}
 \caption{Result of the analysis with different voltage peak values for the pulse. $V_{p}$ is the peak value for the pulse, $t_{0}$ bullet formation time, $V_{0}$ bullet formation tension, $x_{\text{dist}}$ and $y_{\text{dist}}$ are the highest distances reached by the bullet from electrode position, $v_{N}$ is the velocity in x direction reached inside the nozzle, $v_{A}$ is the velocity of propagation in air in x direction.}
 \label{tab:elio_d}
\end{table}
\end{comment}

\begin{table}
 \centering
 \begin{tabular}{cccccc}
 \toprule
 $V_{p}$ [kV]    &$V_{0}$ [kV]    &$x_{\text{dist}}$ [mm]   &$y_{\text{dist}}$ [mm]   &$v_{N}$ [km/s]   &$v_{A}$ [km/s]\\
 \midrule
 \num{5.66(30)}  &\num{5.55(28)}    &\num{15.99(1)} &\num{1.66(2)}  &\num{37.72(4)} &-\\
 \num{6.59(33)}  &\num{5.71(29)}    &\num{26.08(2)} &\num{1.92(4)}  &\num{46.48(2)} &\num{95.16(6)}\\
 \num{7.22(36)}  &\num{6.01(30)}    &\num{34.55(5)} &\num{2.07(1)}  &\num{59.80(3)} &\num{149.47(9)}\\
 \bottomrule
 \end{tabular}
 \caption{Result of the analysis with different voltage peak values for the pulse. $V_{p}$ is the peak value for the pulse, $V_{0}$ bullet formation tension, $x_{\text{dist}}$ and $y_{\text{dist}}$ are the highest distances reached by the bullet from electrode position, $v_{N}$ is the velocity in x direction reached inside the nozzle, $v_{A}$ is the velocity of propagation in air in x direction.}
 \label{tab:elio_d}
\end{table}


\subsection{Flow influence}
As said before, neutral gas velocity is very little in comparision of the ionization front velocity of the bullet, we expect that changes in the gas flow will have an effect on the quantity of helium present in the nozzle more then on the energy of the atoms and electrons. We try four different flows with the same peak voltage, in figure \ref{fig:elio_d} and table \ref{tab:elio_d} are presented the results of the analysis.

Also for this parameter, the discharge starts always in the same voltage interval, around \SI{5.7}{\kilo\volt}. For gas flows of $\num{1}$ and \SI{2}{\liter/\minute} bullet luminosity decreases in a time interval around \SI{100}{\nano\second} once it exits the nozzle, for higher flows it has a luminosity clearly stronger then background for a longer time interval, around \SI{200}{\nano\second} (figure \ref{fig:elio_035} (a)).

From the barycenter coordinates (figure \ref{fig:elio_d} (b) - (c)) we see as more gas flow seems to be correlated with lower mobility of the bullet: we have shorter distances with higher flows. For maximum tried flow, even if bullet luminosity is more then the background, the bullet travels a very little distance from the exit of the nozzle and it's barycenter remains almost fixed in space.

Dimensions of the bullet (figure \ref{fig:elio_d} (d) - (e)) are not very  influenced by gas flow, but we can see little inverse proportionality. For stronger flows we have always littler diameters inside the nozzle, in both directions.

Velocities shows how the bullet behaviour changes also inside the nozzle varying gas flow (figure \ref{fig:elio_d} (f)). With little flow, we find a costant displacement velocity in the nozzle, while with more flow the bullet slows down when moving (as shown in \cite{Jarrige_2010}), with a decrease in velocity of \SI{7.58(4)}{\kilo\meter/\second/\milli\meter}.
Once the bullet exits the nozzle, we see how it almost stops when there is high flow, while it shows the same behaviour if flow is under \SI{3}{\liter/\minute} and it has maximum velocity for the lowest gas flow, maintaining it for longer time.


Can be concluded that with higher flow we have bullet with longer lifetime and littler mobility, reaching lesser distances with lesser velocity.


\begin{figure}
 \centering
 \subfloat[Mean luminosity.]{
    \includegraphics[width=0.48\textwidth]{Images/Shape/elio_035_Im.png}
 }
 \hfill
 \subfloat[x barycenter.]{
    \includegraphics[width=0.48\textwidth]{Images/Shape/elio_035_xb.png}
 }
 \subfloat[y barycenter.]{
    \includegraphics[width=0.48\textwidth]{Images/Shape/elio_035_yb.png}
 }
 \hfill
 \subfloat[x diameter.]{
    \includegraphics[width=0.48\textwidth]{Images/Shape/elio_035_dx.png}
 }
 \subfloat[y diameter.]{
    \includegraphics[width=0.48\textwidth]{Images/Shape/elio_035_dy.png}
 }
 \hfill
 \subfloat[x velocity.]{
    \includegraphics[width=0.48\textwidth]{Images/Shape/elio_035_vx.png}
 }
 \caption{Results of the analysis for different gas flow (opening time of \SI{3.5}{\micro\second} without target).}
 \label{fig:elio_d}
\end{figure}

\begin{comment}
\begin{table}
 \centering
 \begin{tabular}{ccccccc}
 \toprule
 Flow [L/min]    &$t_{0}$ [ns] &$V_{0}$ [kV]    &$x_{\text{dist}}$ [mm]   &$y_{\text{dist}}$ [mm]   &$v_{N}$ [km/s]   &$v_{A}$ [km/s]\\
 \midrule
 \num{1}  &\num{121.9}  &\num{5.87(8)}    &\num{30.65(3)} &\num{2.75(1)}  &\num{54.33(10)} &\num{113.97(9)}\\
 \num{2}  &\num{120.9}  &\num{5.71(29)}    &\num{26.08(2)} &\num{1.92(4)}  &\num{46.48(2)} &\num{95.16(6)}\\
 \num{3}  &\num{124.0}  &\num{5.41(27)}    &\num{19.38(1)} &\num{1.84(1)}  &\num{39.88(4)} &\num{40.96(10)}\\
 \num{4}  &\num{121.3}  &\num{5.90(12)}    &\num{22.80(3)} &\num{2.30(1)}  &\num{58.89(14)} &\num{61.94(44)}\\
 \bottomrule
 \end{tabular}
 \caption{Result of the analysis with different gas flows. $t_{0}$ is the bullet formation time, $V_{0}$ the bullet formation tension, $x_{\text{dist}}$ and $y_{\text{dist}}$ are the highest distances reached by the bullet from electrode position, $v_{N}$ is the velocity in x direction reached inside the nozzle, $v_{A}$ is the velocity of propagation in air in x direction.}
 \label{tab:elio_d}
\end{table} 
\end{comment}

\begin{table}
 \centering
 \begin{tabular}{cccccc}
 \toprule
 Flow [L/min]    &$V_{0}$ [kV]    &$x_{\text{dist}}$ [mm]   &$y_{\text{dist}}$ [mm]   &$v_{N}$ [km/s]   &$v_{A}$ [km/s]\\
 \midrule
 \num{1}  &\num{5.87(8)}    &\num{30.65(3)} &\num{2.75(1)}  &\num{54.33(10)} &\num{113.97(9)}\\
 \num{2}  &\num{5.71(29)}    &\num{26.08(2)} &\num{1.92(4)}  &\num{46.48(2)} &\num{95.16(6)}\\
 \num{3}  &\num{5.41(27)}    &\num{19.38(1)} &\num{1.84(1)}  &\num{39.88(4)} &\num{40.96(10)}\\
 \num{4}  &\num{5.90(12)}    &\num{22.80(3)} &\num{2.30(1)}  &\num{58.89(14)} &\num{61.94(44)}\\
 \bottomrule
 \end{tabular}
 \caption{Result of the analysis with different gas flows. $V_{0}$ the bullet formation tension, $x_{\text{dist}}$ and $y_{\text{dist}}$ are the highest distances reached by the bullet from electrode position, $v_{N}$ is the velocity in x direction reached inside the nozzle, $v_{A}$ is the velocity of propagation in air in x direction.}
 \label{tab:elio_d}
\end{table}


\subsection{Insulator target}
Plasma produced with this source will be applied on biological tissues, so it's fundamental to know how bullet behaviour changes with different targets.

An insulator target is not expected to affect the behaviour until the bullet reaches it. An intersting effect is seen after the impact of the bullet with the target, as in figure \ref{fig:elio_ins}, where the plasma seems to deposit charge on the insulator in a circular shape.
\begin{figure}
 \centering
 \subfloat[Bullet propagation \SI{437.6}{\nano\second}.]{
    \includegraphics[width=0.42\textwidth]{Images/Shape/elio_ins1.png}
 }
 \hfill
 \subfloat[Bullet impact \SI{457.6}{\nano\second}.]{
    \includegraphics[width=0.42\textwidth]{Images/Shape/elio_ins2.png}
 }
 \hfill
 \subfloat[Charge deposition \SI{577.6}{\nano\second}.]{
    \includegraphics[width=0.42\textwidth]{Images/Shape/elio_ins3.png}
 }
 \caption{Impact and charge deposition on an insulator target.}
 \label{fig:elio_ins}
\end{figure}


In figure \ref{fig:elio_c_xb} we can see barycenter's x-coordinate of the bullet and its diameter. While for the lowest tension value we see that the bullet stops before it reaches the target, for the other two sets it clearly stops at the target position. Distance values are comparable with those without target (figure \ref{fig:elio_d}), a little higher for medium tension value, as the bullet instead of stopping at a position before it (as it does whitout target) collide with the target.
Bullet width is exactly the same with or without target during it's propagation. We can see changes only when it impacts on the target and shrinks rapidly.
Displacement velocities inside the nozzle and propagation velocities in air are of the same order of values without target, as we can see in figure \ref{fig:elio_c_vx}. We find higher $v_{A}$ only for medium tension value, where it reaches \SI{118.98(7)}{\kilo\meter/\second}, due to the rapid collision with the target instead of slowing down and stopping.
\begin{figure}
 \centering
 \subfloat[Barycenter position.]{
    \includegraphics[width=0.48\textwidth]{Images/Shape/elio_c_xb.png}
 }
 \hfill
 \subfloat[Bullet diameter.]{
    \includegraphics[width=0.48\textwidth]{Images/Shape/elio_c_dx.png}
 }
 \caption{Barycenter motion and diameter in x direction, for different voltage peak values.}
 \label{fig:elio_c_xb}
\end{figure}

\paragraph{Charge deposition}
Once the bullet reaches the target we can see how the glowing region enlarges on its surface from the plot of y diameter, as in figure \ref{fig:elio_c_ylim}. To evaluate the expansion diameter and velocity, we can plot the differance between y limit and y barycenter once the bullet reaches the target. In figure \ref{fig:elio_c_ylim} we can see that the diameter of the maximum expansion is little higher for larger tension value, as well as it's expansion velocity: for opening time of \SI{3.5}{\micro\second} we find a diameter of \SI{6.54(4)}{\milli\meter} and an expansion velocity of \SI{31.36(25)}{\kilo\meter/\second}, while for \SI{4.0}{\micro\second} we find \SI{7.40(2)}{\milli\meter} and \SI{64.04(39)}{\kilo\meter/\second}. It is an expected result as the bullet have more energy with an higher peak value, the velocity along x is larger and we have also higher expansion velocity.

\begin{figure}
 \centering
 \subfloat[Bullet diameter.]{
    \includegraphics[width=0.48\textwidth]{Images/Shape/elio_c_dy.png}
 }
 \hfill
 \subfloat[Bullet expansion in time.]{
    \includegraphics[width=0.48\textwidth]{Images/Shape/elio_c_ylim.png}
 }
 \caption{Bullet expansion in the y direction. We see the absolute value for the diameter for all barycenter positions (a) and a zoom on the expansion for times after it reaches the target (b), where y contour values are calculated as differance with the barycenter y coordinate. Linear fits show how much the diameter enlarges in time.}
 \label{fig:elio_c_ylim}
\end{figure}


\subsection{Distant conductive target}
In measure set A the target is positioned at \SI{24}{\milli\meter} from the bench, correponding to \SI{30}{\milli\meter} from the electrode. A conductive target is expected to affect bullet behaviour as it reachs its proximity, as there are free charges on the surface. If with an insulator we see charge deposition, with a conductor we see a rapid backstream (with a motion inverse to that of the bullet) with great luminosity when the bullet impacts on it, as in figure \ref{fig:elio_met}.
\begin{figure}
 \centering
 \subfloat[Bullet propagation \SI{371.2}{\nano\second}.]{
    \includegraphics[width=0.42\textwidth]{Images/Shape/elio_met1.png}
 }
 \hfill
 \subfloat[Bullet impact \SI{391.2}{\nano\second}.]{
    \includegraphics[width=0.42\textwidth]{Images/Shape/elio_met2.png}
 }
 \hfill
 \subfloat[Backstream \SI{411.2}{\nano\second}.]{
    \includegraphics[width=0.42\textwidth]{Images/Shape/elio_met3.png}
 }
 \caption{Impact and backstream with a copper target.}
 \label{fig:elio_met}
\end{figure}

\paragraph{Bullet dynamics}
Considering only the bullet, separated from the backstream, first thing to note with this target, is how all the system is more reactive to tension values: as we can see in figure \ref{fig:elio_a_Im}, we find that plasma bullet luminosity is enanched in correspondence of the second tension peak (negative), even if it doesn't move from the target as shown in figure \ref{fig:elio_a_xb} (a). Thanks to this second peak, bullet lifetime is longer, it shows luminosity even after \SI{900}{\nano\second} from the start.
\begin{figure}
 \centering
 \subfloat[Voltage waveforms.]{
    \includegraphics[width=0.43\textwidth]{Images/Shape/elio_a_V.png}
 }
 \hfill
 \subfloat[Bullet luminosity.]{
    \includegraphics[width=0.43\textwidth]{Images/Shape/elio_a_Im.png}
 }
 \caption{Voltage applied to the elctrode and bullet luminosity for different measure sets. Note the second luminosity peak in correspondence of the negative voltage peak.}
 \label{fig:elio_a_Im}
\end{figure}


Barycenter motion shows also some peculiarities: even with low tension the bullet reaches the target, showing larger propagation velocity then with other configurations. It's diameter doesn't shows different behaviour, it stays constant until nozzle exit and remains constant until it reaches the target and shrinks to a point. As said before, velocities of propagation in air are generally larger then whitout target, they arrives from \SI{112.04(5)}{\kilo\meter/\second} for low voltage peak to \SI{160.55(9)}{\kilo\meter/\second} for high voltage peak.
\begin{figure}
 \centering
 \subfloat[Barycenter.]{
    \includegraphics[width=0.43\textwidth]{Images/Shape/elio_a_xb.png}
 }
 \hfill
 \subfloat[Diameter.]{
    \includegraphics[width=0.43\textwidth]{Images/Shape/elio_a_dx.png}
 }
 \subfloat[Velocity.]{
    \includegraphics[width=0.43\textwidth]{Images/Shape/elio_a_vx.png}
 }
 \caption{Bullet barycenter coordinate, diameter and velocity along x direction, with copper target. When the bullet reaches the target, it shrinks and stops.}
 \label{fig:elio_a_xb}
\end{figure}


\paragraph{Current measure}
Thanks to the conductive target, we can measure current intensity flowing in it due to charge transported by bullets, and, more intersting, the relation between impact time and when it's measured peak current.
Current intensity it's shown in figure \ref{fig:elio_a_icurr}, we can see a positive peak, more influenced by tension peak value, and a negative peak in correspondence of the negative tension peak and of the second luminosity peak for the bullet. It's important to note that the time interval where we have high luminosity for the bullet and its position is around the target position, is the same time interval where we measure current.
\begin{figure}
 \centering
 \includegraphics[width=0.6\textwidth]{Images/Shape/elio_a_icurr.png}
 \caption{Current intensity measured in set A, with different voltage peak values.}
 \label{fig:elio_a_icurr}
\end{figure}

In the graphs we can distinguish a starting time for the peak current, and a time where it reaches it's maximum value. Even if it's difficult to establish precisely the starting time, we can see as both of them are larger for lower tension values, we have lower and slower current intensity when bullets have less energy. To study the relation between those values, we extrapolate relevant times and the bullet distance from the target when we start to see current, presented in table \ref{tab:elio_a_times}. From the values we can see that we begin to measure current before the bullet barycenter reaches the target, and even maximum current is measured before the impact time of the bullet for medium and high tension value.
We can estimate a distance of interaction as the differance between bullet barycenter position and target position. This distance becames larger when we set an higher voltage for the pulse, when the bullet velocity is higher.

Current measure can be explained if the bullet motion is associated with charge motion: as the bullet travels towards the target, charges with highest speed reach the target before the barycenter and we begin to measure current before it reaches the target. Once the bullet loses luminosity, also current stops to flow in the target. When the bullet has higher energy, charges have more speed, arrive before to the target and we will lose lesser charges during propagation, so we will measure more current, with a peak that starts earlyer.
\begin{table}
 \centering
 \begin{tabular}{cccccc}
  \toprule
  $V_{p}$ [kV]  &$i_{p}$ [mA]   &$t_{0,i}$ [ns] &$t_{p,i}$ [ns] &$t_{\text{imp}}$ [ns]  &$x_{\text{target}} - x(t_{0,i})$ [mm]\\
  \midrule
  \num{5.7(3)}  &\num{17.56(88)}    &\num{459.1(150)}   &\num{509.2(150)}   &\num{499.1(150)}   &\num{7.20(1)}\\
  \num{6.6(3)}  &\num{32.63(163)}    &\num{391.6(150)}   &\num{451.6(150)}   &\num{511.6(150)}   &\num{9.22(1)}\\
  \num{7.2(4)}  &\num{41.50(62)}    &\num{350.3(150)}   &\num{410.3(150)}   &\num{450.3(150)}   &\num{15.18(2)}\\
  \bottomrule
 \end{tabular}
 \caption{Values extrapolated from current measure and bullet barycenter motion for set A. $V_{p}$ is the voltage peak for the pulse, $i_{p}$ is the current peak value, $t_{0,i}$ is the starting time for the current value, $t_{p,i}$ is the time of the peak value, $t_{\text{imp}}$ is the time of the impact of bullet on target, $x_{\text{target}} - x(t_{0,i})$ is the distance between bullet position when we start to measure current and target position.}
 \label{fig:elio_a_times}
\end{table}


\paragraph{Backstream}
Dynamics of the backstream from the target is difficult to analyze as it's a rapid phenomenon with too high luminosity. We see that the frames right after the impact of the bullet on target, present another zone with high luminosity right after nozzle exit. This other glow has stable position and diameter, as can be seen in figure \ref{fig:elio_a_back}, with a slow tendency to shift towards the target when voltage is positive and towards the electrode when it's negative.
The phenomenon can be explained as a very rapid charge stream from the target to the electrode, that stops when meets the nozzle, but to study it are necessary other specific measures with lower velocities (for example increasing gas flow, see before) or more time resolution.
\begin{figure}
 \centering
 \subfloat[Mean luminosity.]{
    \includegraphics[width=0.43\textwidth]{Images/Shape/elio_a_Im_back.png}
 }
 \hfill
 \subfloat[Barycenter along x.]{
    \includegraphics[width=0.43\textwidth]{Images/Shape/elio_a_xb_back.png}
 }
 
 \subfloat[Diameter along x.]{
    \includegraphics[width=0.43\textwidth]{Images/Shape/elio_a_dx_back.png}
 }
 \hfill
 \subfloat[Diameter along y.]{
    \includegraphics[width=0.43\textwidth]{Images/Shape/elio_a_dy_back.png}
 }
 \caption{Backstream features versus time, measure set A.}
 \label{fig:elio_a_back}
\end{figure}


\subsection{Close conductive target}
In measure set B the target is positioned at \SI{32}{\milli\meter} from the bench, correponding to \SI{22}{\milli\meter} from the electrode.
With a conductive target this close it's more difficult to see the bullet evolution, as it is more rapid and the backstream has even higher relevance.

In figure \ref{fig:elio_b_dyn} are presented the results of the analysis. For all the sets we see always high luminosity, even after a time interval corresponding to the end of the voltage peak. As we can see from barycenter and diameter values, the peculiarity of this set is that for higher voltage peak values, we cannot even see the bullet, as it closes on the target in a sort of large plume right after nozzle exit. The other two sets shows the same behaviour described before, with even higher velocities: for medium voltage we find a velocity of \SI{59.30(5)}{\kilo\meter/\second} inside the nozzle and a velocity of \SI{174.44(8)}{\kilo\meter/\second} in air.
\begin{figure}
 \centering
 \subfloat[Mean luminosity.]{
    \includegraphics[width=0.43\textwidth]{Images/Shape/elio_b_Im.png}
 }
 \hfill
 \subfloat[Barycenter along x.]{
    \includegraphics[width=0.43\textwidth]{Images/Shape/elio_b_xb.png}
 }
 
 \subfloat[Diameter along x.]{
    \includegraphics[width=0.43\textwidth]{Images/Shape/elio_b_dx.png}
 }
 \hfill
 \subfloat[Diameter along y.]{
    \includegraphics[width=0.43\textwidth]{Images/Shape/elio_b_vx.png}
 }
 \caption{Analysis of measure set B.}
 \label{fig:elio_b_dyn}
\end{figure}

\paragraph{Current measure}
As done with set A, for low and medium voltage it's possible to analyze time and space intervals between bullet barycenter motion and current measure. Obtained values are shown in table \ref{tab:elio_b_times}.
\begin{table}
 \centering
 \begin{tabular}{cccccc}
  \toprule
  $V_{p}$ [kV]  &$i_{p}$ [mA]   &$t_{0,i}$ [ns] &$t_{p,i}$ [ns] &$t_{\text{imp}}$ [ns]  &$x_{\text{target}} - x(t_{0,i})$ [mm]\\
  \midrule
  \num{5.7(3)}  &\num{47.41(237)}    &\num{370.6(150)}   &\num{450.6(150)}   &\num{450.6(150)}   &\num{11.34(2)}\\
  \num{6.6(4)}  &\num{77.41(387)}    &\num{311.2(150)}   &\num{411.2(150)}   &\num{391.2(150)}   &\num{13.61(4)}\\
  \bottomrule
 \end{tabular}
 \caption{Values extrapolated from current measure and bullet barycenter motion for set B. $V_{p}$ is the voltage peak for the pulse, $i_{p}$ is the current peak value, $t_{0,i}$ is the starting time for the current value, $t_{p,i}$ is the time of the peak value, $t_{\text{imp}}$ is the time of the impact of bullet on target, $x_{\text{target}} - x(t_{0,i})$ is the distance between bullet position when we start to measure current and target position.}
 \label{fig:elio_a_times}
\end{table}

Again we find that current flows in target before the impact of the bullet, when it is distant from the target (parameter $x_{\text{target}} - x(t_{0,i})$ in table), but current rises more rapidly if the bullet has more energy. Comparing these values with results for set A, we find that for the same voltage peak value, the current peak starts faster and reaches an higher peak value. Again, hypotizing that the bullet motion is correlated with charge motion, an earlier impact point (as in set B compared to set A) means that we will lose lesser charge during the propagation and measure an higher current value on target.


\section{Neon flow}
The second noble gas used to produce cold plasma is neon, the easyier one to ignite
and more luminous between used gas mixture of noble elements with low weight.
Neon it's an element with standard atomic weight of $\num{20.180}$, 1st ionization energy of \SI{21.565}{\electronvolt} and 2nd ionization energy of \SI{40.963}{\electronvolt} and several lines with wavelenght in visible spectrum. It's easy to ignite because when the voltage pulse is applied to the electrode, even for low frequencies, the gas presents high luminosity in a wide space. This emissivity it's the reason why neon is commonly used to build neon lamps.

Measures with this gas are done to compare bullets with entirely different gas features and equivalent setups. As plasma emission begins with lower voltage and frequencies for electrode pulses compared to helium, we set a low working frequency $f = \SI{1}{\kilo\hertz}$ and three opening times nearly equivalent to low, medium and high voltage from before. The different measures setup are presented in table \ref{tab:setups}.

\subsection{Free bullet}
We want to see plasma dynamics whitout a target, changing between low, medium and high tension peak value. Results are shown in figure \ref{fig:neon_d} and table \ref{tab:neon_d}.
First thing to notice is that plasma is produced not specifically at electrode's end where the electric field should be higher, as in helium, but a little before, with lower luminosity. Luminosity evolution from there is similar to what we have in helium, but there is a long lifed residual luminosity that lasts longer.
Studying barycenter motion and it's diameter, we see that until the end of the nozzle we have similar dynamics, after it, we can point out two main differances:
\begin{itemize}
 \item bullets come out of the nozzle more difficultly and lose velocity more rapidly, covering littler distances. Propagation velocity in air rises exiting the nozzle (until \SI{160.32(7)}{\kilo\meter/\second} with high voltage), but goes down abruptily and the bullet slows down, as its luminosity lowers.
 \item differently from helium, bullet's width rise a little increasing voltage. As we can see from figure, bullet's diameter is almost a constant value during the motion, increasing only before nozzle exit, however we have always higher values with higher voltage, suggesting that the bullet is more large with higher voltage.
\end{itemize}

This behaviour can be explained as the bullet is the result of more excitation reactions that produces emission with wavelenght in visible spectrum, and those reactions have lower threshold energy. If reactions needs less energetic electrons, we have more of them that partecipates, spreading more the glow, and more of them that loses energy once they are in contact with air outside the nozzle, distant from the high voltage electrode.

\begin{figure}
 \centering
 \subfloat[Voltage waveform.]{
    \includegraphics[width=0.46\textwidth]{Images/Shape/neon_a_V.png}
 }
 \hfill
 \subfloat[Mean luminosity.]{
    \includegraphics[width=0.46\textwidth]{Images/Shape/neon_d_Im.png}
 }
 
 \subfloat[x barycenter.]{
    \includegraphics[width=0.46\textwidth]{Images/Shape/neon_d_xb.png}
 }
 \hfill
 \subfloat[x velocity.]{
    \includegraphics[width=0.46\textwidth]{Images/Shape/neon_d_vx.png}
 }
 
 \subfloat[x diameter.]{
    \includegraphics[width=0.46\textwidth]{Images/Shape/neon_d_dx.png}
 }
 \hfill
 \subfloat[y diameter.]{
    \includegraphics[width=0.46\textwidth]{Images/Shape/neon_d_dy.png}
 }
 \caption{Results of the analysis for setup D wiith neon gas, for different peak voltage values.}
 \label{fig:neon_d}
\end{figure}


\begin{table}
 \centering
 \begin{tabular}{ccccc}
 \toprule
 $V_{p}$ [kV]    &$V_{0}$ [kV]    &$x_{\text{dist}}$ [mm]   &$v_{N}$ [km/s]   &$v_{A}$ [km/s]\\
 \midrule
 \num{4.81(24)}  &\num{4.81(24)}    &\num{18.13(5)} &\num{34.59(3)} &\num{43.19(7)}\\
 \num{5.33(27)}  &\num{5.20(26)}    &\num{22.42(6)} &\num{42.82(4)} &\num{61.17(5)}\\
 \num{6.05(30)}  &\num{5.16(26)}    &\num{26.60(9)} &\num{62.92(5)} &\num{160.32(7)}\\
 \bottomrule
 \end{tabular}
 \caption{Result of the analysis with different voltage peak values for the pulse. $V_{p}$ is the peak value for the pulse, $V_{0}$ bullet formation tension, $x_{\text{dist}}$ is the highest distances reached by the bullet from electrode position, $v_{N}$ is the velocity in x direction reached inside the nozzle, $v_{A}$ is the velocity of propagation in air.}
 \label{tab:neon_d}
\end{table}


For both helium and neon inside the nozzle there is bullet propagation, it is possible to compare how the velocity $v_{N}$ changes varying voltage peak value for the gasses. In figure \ref{fig:hene_d_vn} there is the plot for the two gasses, as we can see with few points it's not possible to extrapolate accurately the behavior, with a linear plot we find that helium bullet's velocity grows with a slope of \SI{14.55(240)}{\kilo\meter/\second \kilo\volt}, while for neon bullet's velocity we find \SI{23.23(366)}{\kilo\meter/\second \kilo\volt}. Those values are of the same magnitude, but not so compatible, suggesting that bullet's velocity may depend from gas composition not only with a different threshold value.
\begin{figure}
 \centering
 \includegraphics[width=0.52\textwidth]{Images/Shape/hene_vN.png}
 \caption{Plot of bullet's velocity inside the nozzle increasing voltage, for helium and neon.}
 \label{fig:hene_d_vn}
\end{figure}


\subsection{Insulator target}
Also for this gas we see as an insulator target, positioned at \SI{30}{\milli\meter} from the electrode, influences bullet dynamics.
Results are similar with helium or neon bullets: palsma forms near the electrode, exits from the nozzle and the impact on the target produce a round shaped figure on it, as in figure \ref{fig:elio_ins}. We can see also in neon that, for low voltage peak value, the bullet doesn't reaches the target, as the maximum distance travelled is lower then target distance, seen with setup D. This suggests that the target does not influences the bullet when it is not in its proximity.
A peculiarity with neon is that we can see a second, very low, peak in luminosity after long time (measure taken only with medium voltage in figure), corresponding to the second voltage peak of the pulse. Once again we see that neon is more reactive then helium, were this peak doesn't shows.
\begin{figure}
 \centering
 \subfloat[Mean luminosity.]{
    \includegraphics[width=0.46\textwidth]{Images/Shape/neon_c_dx.png}
 }
 \hfill
 \subfloat[Barycenter position.]{
    \includegraphics[width=0.46\textwidth]{Images/Shape/neon_c_xb.png}
 }
 \caption{Mean luminosity and barycenter motion for neon gas in setup C.}
 \label{fig:neon_c_xb}
\end{figure}


The circle produced by bullet impact on target can be analyzed as done before, showing how the figure expands in y direction. From figure \ref{fig:neon_c_ylim} we can extrapolate a velocity of \SI{5.01(29)}{\kilo\meter/\second} for medium voltage and \SI{30.76(25)}{\kilo\meter/\second} for high voltage. They are both lower compared to helium, as expected if we consider that the bullet slows down more rapidly once it's outside the nozzle.
\begin{figure}
 \centering
 \includegraphics[width=0.46\textwidth]{Images/Shape/neon_c_ylim.png}
 \caption{Neon plasma expansion on insulator target, for different voltage peak values.}
 \label{fig:neon_c_ylim}
\end{figure}


\subsection{Conductive target}
As with helium, a conductive target has an influence in all neon bullet's dynamics, in particular we see a more rapid passage from the bullet shape to the elongated one that gives rise to the backstream described in figure \ref{fig:elio_met}.

In figure \ref{fig:neon_ab_xb} we see that the distant target (setup A) is reached by the bullet with medium and high voltage, but there isn't a change in luminosity. Near nozzle exit bullets speed up, reach rapidly the target, stretch and then reduces their diameter to the impact point.
For a close target we see a change in luminosity in correspondence of the negative peak and of the second positive peak. With high voltage we loose completely the bullet shape once plasma exits the nozzle, all the space between nozzle exit and target is completely filled by glowing gas.
\begin{figure}
 \centering
 \subfloat[Mean luminosity, $\Delta x_{\text{targ}} = \SI{30}{\milli\meter}$.]{
    \includegraphics[width=0.46\textwidth]{Images/Shape/neon_a_Im.png}
 }
 \hfill
 \subfloat[Mean luminosity, $\Delta x_{\text{targ}} = \SI{22}{\milli\meter}$.]{
    \includegraphics[width=0.46\textwidth]{Images/Shape/neon_b_Im.png}
 }
 
 \subfloat[x barycenter, $\Delta x_{\text{targ}} = \SI{30}{\milli\meter}$]{
    \includegraphics[width=0.46\textwidth]{Images/Shape/neon_a_xb.png}
 }
 \hfill
 \subfloat[x barycenter, $\Delta x_{\text{targ}} = \SI{22}{\milli\meter}$]{
    \includegraphics[width=0.46\textwidth]{Images/Shape/neon_b_xb.png}
 }
 
 \subfloat[x diameter, $\Delta x_{\text{targ}} = \SI{30}{\milli\meter}$]{
    \includegraphics[width=0.46\textwidth]{Images/Shape/neon_a_dx.png}
 }
 \hfill
 \subfloat[x diameter, $\Delta x_{\text{targ}} = \SI{22}{\milli\meter}$]{
    \includegraphics[width=0.46\textwidth]{Images/Shape/neon_b_dx.png}
 }
 \caption{Results of the analysis for setup A and B with neon gas, for different peak voltage values.}
 \label{fig:neon_ab_xb}
\end{figure}


\paragraph{Current measure}
Current measures are different respect to measure done with helium gas, here results suggest a relation between luminosity intensity and current intensity. As we can see in figure \ref{fig:neon_ab_icurr} ad in table \ref{tab:neon_ab_ival}, a current peak is measured only for high voltage with distant target and for medium and high voltage with close target. If emission intensity is too low when the bullet enter in the range of interaction with the target, current intensity is not measured, with higher voltage the bullet is more rapid, reaches the target before the loss in luminosity and we measure current.

Starting time of current peak is again precedent to the impact time of the bullet with the target. Interaction distances presented in table are similar to the ones found for helium.
\begin{figure}
 \centering
 \subfloat[$\Delta x_{\text{targ}} = \SI{30}{\milli\meter}$.]{
    \includegraphics[width=0.46\textwidth]{Images/Shape/neon_a_i.png}
 }
 \hfill
 \subfloat[$\Delta x_{\text{targ}} = \SI{22}{\milli\meter}$.]{
    \includegraphics[width=0.46\textwidth]{Images/Shape/neon_b_i.png}
 }
 \caption{Current intensity measured in set A and B, with different voltage peak values.}
 \label{fig:neon_ab_icurr}
\end{figure}

\begin{table}
 \centering
 \begin{tabular}{ccccccc}
  \toprule
  $\Delta x_{\text{targ}}$ [mm] &$V_{p}$ [kV]  &$i_{p}$ [mA]   &$t_{0,i}$ [ns] &$t_{p,i}$ [ns] &$t_{\text{imp}}$ [ns]  &$x_{\text{target}} - x(t_{0,i})$ [mm]\\
  \midrule
  \num{30}  &\num{6.05(30)}  &\num{6.87(34)}    &\num{329(15)}   &\num{477(15)}   &\num{453(21)}   &\num{19.96(3)}\\
  \midrule
  \multirow{2}*{\num{22}}   &\num{5.33(27)}  &\num{7.72(39)}    &\num{347(15)}   &\num{447(15)}   &\num{447(15)}   &\num{13.06(3)}\\
                            &\num{6.05(30)}  &\num{14.56(73)}    &\num{300(15)}   &\num{398(15)}   &\num{380(21)}   &\num{13.66(4)}\\
  \bottomrule
 \end{tabular}
 \caption{Values extrapolated from current measure and bullet barycenter motion for neon, set A and B. $\Delta x_{\text{targ}}$ is the distance between electrode and target, $V_{p}$ is the voltage peak for the pulse, $i_{p}$ is the current peak value, $t_{0,i}$ is the starting time for the current value, $t_{p,i}$ is the time of the peak value, $t_{\text{imp}}$ is the time of the impact of bullet on target, $x_{\text{target}} - x(t_{0,i})$ is the distance between bullet position when we start to measure current and target position.}
 \label{fig:neon_ab_ival}
\end{table}


\paragraph{Backstream}
Backstream phenomenon in neon is easily observable, where the bullet reaches the target we find high luminosity around nozzle exit. It presents itself as zone of glowing gas with large width and near constant height, with a barycenter motion that shift direction towards the electrode or towards the target, with so little motion that is not possible to infer it's dynamics. Once the pulse voltage lowers, also the backstream fade away, as the bullet that impacts on target. 

\begin{figure}
 \centering
 \subfloat[Mean luminosity, $\Delta x_{\text{targ}} = \SI{30}{\milli\meter}$.]{
    \includegraphics[width=0.43\textwidth]{Images/Shape/neon_a_back_Im.png}
 }
 \hfill
 \subfloat[Mean luminosity,  $\Delta x_{\text{targ}} = \SI{22}{\milli\meter}$.]{
    \includegraphics[width=0.43\textwidth]{Images/Shape/neon_b_back_Im.png}
 }
 
 \subfloat[Barycenter along x, $\Delta x_{\text{targ}} = \SI{30}{\milli\meter}$.]{
    \includegraphics[width=0.43\textwidth]{Images/Shape/neon_a_back_xb.png}
 }
 \hfill
 \subfloat[Barycenter along x,  $\Delta x_{\text{targ}} = \SI{22}{\milli\meter}$.]{
    \includegraphics[width=0.43\textwidth]{Images/Shape/neon_b_back_xb.png}
 }
 
 \subfloat[Diameter along x, $\Delta x_{\text{targ}} = \SI{30}{\milli\meter}$.]{
    \includegraphics[width=0.43\textwidth]{Images/Shape/neon_a_back_dx.png}
 }
 \hfill
 \subfloat[Diameter along x,  $\Delta x_{\text{targ}} = \SI{22}{\milli\meter}$.]{
    \includegraphics[width=0.43\textwidth]{Images/Shape/neon_b_back_dx.png}
 }
 
 %\subfloat[Diameter along y, $\Delta x_{\text{targ}} = \SI{30}{\milli\meter}$.]{
 %   \includegraphics[width=0.43\textwidth]{Images/Shape/neon_a_back_dy.png}
 %}
 %\hfill
 %\subfloat[Diameter along y.]{
 %   \includegraphics[width=0.43\textwidth]{Images/Shape/neon_b_back_dy.png}
 %}
 \caption{Analysis of the backstream in measure set A and B for neon.}
 \label{fig:neon_ab_tot_back}
\end{figure}


\section{Argon flow}
The third noble gas used it's Argon, it is the harder one to produce plasma with at atmosferic pressure. Argon has a standard atomic weight of $\num{39.948}$, 1st ionization energy of \SI{15.760}{\electronvolt} and 2nd ionization energy of \SI{27.629}{\electronvolt}.
The peculiarity of this gas is its tendency to form filaments of glowing plasma, long in the electrode axis direction and short in height, instead of a large zone of plasma column as seen with helium or neon. Figure \ref{fig:argonex} shows an example where there is a grounded ring made of copper positioned around the outer nozzle, that hallows to produce plasma with lower voltage and repetition rate. For our measurements we use a frequency and an opening time for the voltage pulse that hallow to observe plasma with and whitout this ring, to see differances between the setups.
\begin{figure}
 \centering
 \includegraphics[width=0.6\textwidth]{Images/Shape/argon_b_120.png}
 \caption{Example of discharge with argon in setup B, after \SI{120}{\nano\second} from discharge start. Note the grounded ring that covers a portion of the nozzle after the electrode.}
 \label{fig:argonex}
\end{figure}


Along the production of filaments, argon shows also the tendency to transite from a glow to an arc discharge when a target is positioned near it, as in figure \ref{fig:arcex}. As we remembered in chapter \ref{ch:electric}, when there is a plasma arc there is also large current intensity circulating in plasma and in the target, a condition that we want to avoid for blood coagulation with cold plasma. For this reason we use the target at one single position, one where we see plasma interaction with the target whitout arc transition.
\begin{figure}
 \centering
 \includegraphics[width=0.6\textwidth]{Images/Shape/arcex.png}
 \caption{Example of arc discharge with argon and a target positioned \SI{22}{\milli\meter} from the electrode, after \SI{500}{\nano\second} from discharge start. The arc covers all the gap between electrode on the left and conductive target on the right.}
 \label{fig:arcex}
\end{figure}

\subsection{Dynamic example}
We take measurements in setup B as example: argon flow of \SI{2}{\liter/\min}, presence of the grounded ring outside the nozzle and absence of target.
Formation, expulsion and propagation of argon plasma is phenomenologically different if compared to helium or neon plasma, but we can always separate two phases: inside the nozzle and outside it.

\paragraph{Plasma formation}
Inside the nozzle we observe plasma formation as filaments that start on the electrode and reach nozzle's walls, with different lenghts for each frame. There isn't a uniform propagation of those filaments toward the exit, however we see that after a time of \SI{180}{\nano\second}, every filament covers all the distance between electrode and nozzle exit (\SI{10}{\milli\meter}), as in figure \ref{fig:arnoz_prop}. Hypotizing the definition of a propagating front composed by the filament's ends it will have a velocity of, at least, $v_{N} = \SI{55.56}{\kilo\meter/\second}$. 
\begin{figure}
 \centering
 \subfloat[t = \SI{0}{\nano\second}]{
    \includegraphics[width=0.5\textwidth]{Images/Shape/argon_b_0.png}
 }
 
 \subfloat[t = \SI{80}{\nano\second}]{
    \includegraphics[width=0.5\textwidth]{Images/Shape/argon_b_80.png}
 }
 
 \subfloat[t = \SI{180}{\nano\second}]{
    \includegraphics[width=0.5\textwidth]{Images/Shape/argon_b_180.png}
 }
 \caption{Propagation of argon plasma filaments inside the nozzle, measurements setup B.}
 \label{fig:arnoz_prop}
\end{figure}

\paragraph{Plasma expulsion}
After the time interval where we can see plasma only inside the nozzle, it's possible to observe zones with high luminosity also in air. Plasma is expelled as tiny round shaped formations, separated from each other, as in figure \ref{fig:arair_prop}.
\begin{figure}
 \centering
 \includegraphics[width=0.6\textwidth]{Images/Shape/argon_b_220.png}
 \caption{Expulsion of argon plasma in air at \SI{220}{\nano\second} from discharge start, measurements set B.}
 \label{fig:arair_prop}
\end{figure}

For every frame we observe several formations with different positions and diameters, it's possible to study the dynamic of plasma propagation studying the evolution of those parameters.
For each frame we identify every formation that presents luminosity over a limit value, isolating it's shape. From there it's possible to count how many formations there are, what are their mean luminosities, barycenter's coordinates and diameters, as we did with the bullets in helium and neon.
Given a frame at a specific time, we fill histograms for those parameters and extrapolate their distribution. For different times, it's possible to compare those distributions, as it's done for barycenter's coordinate x in figure \ref{fig:argon_xb_evol}. 
\begin{figure}
 \centering
 \includegraphics[width=0.6\textwidth]{Images/Shape/argon_b_xbevol.png}
 \caption{Distribution of barycenter's coordinate x, for plasma formations expelled from the source, at three different times, in measurements setup B. For every time, we see how many formations are observed centered at a specific position $x_b$.}
 \label{fig:argon_xb_evol}
\end{figure}


Given those distributions for every parameter, it's possible to evaluate the mean value for luminosity, barycenter coordinates and diameters at any given time, reconstructing an analysis similar to that presented for other gasses. Results are presented in figure \ref{fig:argon_b}, where the points are the mean value of the distribution and the error is associated with distribution widths.
We see that the number of plasma formations follows voltage waveform: we count more plasma formations right after the voltage positive peak value and right after the voltage negative peak value. Also mean luminosity shows this behavior: we see emission with more intensity right after the positive and the negative voltage peaks, with higher errors during the negative peak, showing that values are more spread around the mean value.
Barycenter motion in the x direction shows that the distribution propagates during the time interval where voltage is positive and stops when it becames zero, at a certain distance from the end of the nozzle. It's relevant to point out that we are studying this propagation only outside the nozzle, so position values starts around \SI{85}{\milli\meter}, that is the end of the nozzle. During the propagation phase it's possible to estimate a velocity for the mean value of those formations with a fit as shown in figure, resulting in $v_A = \SI{30.16(356)}{\kilo\meter/\second}$.
Diameters values are presented to show dimensions of plasma formations. Diameter along x and y direction varies in a range between \num{0.3} and \SI{1.3}{\milli\meter}, with a stretch in the x direction where voltage's absolute value is higher, and a stable value in the y direction around \SI{0.7(2)}{\milli\meter}. 

\begin{figure}
 \centering
 \subfloat[Voltage waveform.]{
    \includegraphics[width=0.48\textwidth]{Images/Shape/argon_b_V.png}
 }
 \hfill
 \subfloat[Number of formations.]{
    \includegraphics[width=0.48\textwidth]{Images/Shape/argon_b_N.png}
 }
 
 \subfloat[Mean luminosity.]{
    \includegraphics[width=0.48\textwidth]{Images/Shape/argon_b_Im.png}
 }
 \hfill
 \subfloat[Barycenter along x.]{
    \includegraphics[width=0.48\textwidth]{Images/Shape/argon_b_xb.png}
 }
 
 \subfloat[Diameter along x.]{
    \includegraphics[width=0.48\textwidth]{Images/Shape/argon_b_dx.png}
 }
 \hfill
 \subfloat[Diameter along y.]{
    \includegraphics[width=0.48\textwidth]{Images/Shape/argon_b_dy.png}
 }
 \caption{Analysis of argon plasma formations, measurements setup B. Every point identifies the mean value of considered parameter between all plasma formations observed at a given time, while the error is associated with the distribution width.}
 \label{fig:argon_b}
\end{figure}


\subsection{Absence of target}
The grounded ring positioned around the nozzle for measurements sets A, B and D, phenomenologically, helps the formation of plasma, hallowing to produce it with lower amplitude and repetition rate of voltage pulses. We try visualize the effects of this ring on plasma dynamics comparing the expulsion of plasma in setups with or without it, respectively measurements set A and F.

Analysis is done as described before for set B, results are in figure \ref{fig:argon_af}, where the voltage waveform is identical to the one presented in figure \ref{fig:argon_b}.
As in set B, we find a peak in plasma formation numbers and luminosity in correspondence of voltage peak values, positive and negative. Intersting thing to note it's that the number of formations is always lower for measurements without the ring, implying that it's expelled more plasma when there is the ring.
Baricenter position changes uniformally for both sets, resulting in a velocity of the mean value of $v_A = \SI{47.90(475)}{\kilo\meter/\second}$ with the ring and $v_A = \SI{39.03(295)}{\kilo\meter/\second}$ without it.
Diameters have an average value higher for set A in both direction, but data is spread a lot for all measurements sets, so it's not possible to note differances outside the error.

Argon plasma dynamics doesn't change behavior introducing the grounded ring on the nozzle, but plasma formations that are expelled reach higher velocities.

\begin{figure}
 \centering
 \subfloat[Number of formations.]{
    \includegraphics[width=0.48\textwidth]{Images/Shape/argon_af_N.png}
 }
 \hfill
 \subfloat[Mean luminosity.]{
    \includegraphics[width=0.48\textwidth]{Images/Shape/argon_af_Im.png}
 }
 
 \subfloat[Barycenter along x.]{
    \includegraphics[width=0.48\textwidth]{Images/Shape/argon_af_xb.png}
 }
 
 \subfloat[Diameter along x.]{
    \includegraphics[width=0.48\textwidth]{Images/Shape/argon_af_dx.png}
 }
 \hfill
 \subfloat[Diameter along y.]{
    \includegraphics[width=0.48\textwidth]{Images/Shape/argon_af_dy.png}
 }
 \caption{Analysis of argon plasma formations, with grounded ring positioned outside the nozzle (set A) and without it (set F).}
 \label{fig:argon_af}
\end{figure}

\subsection{Influence of target}
A target positioned after nozzle exit induces easily an arc transition with argon plasma, to study the influence of the target in plasma dynamics it is necessary to find a distance from the electrode for wich there is not an arc transition. In measurements sets D and E we study plasma formation and evolution with a copper target (the same described before) positioned at \SI{30}{\milli\meter}, respectively with and without grounded ring around the nozzle.


Analysis is done as described before, results are in figure \ref{fig:argon_de}, where the voltage waveform is identical to the one in figure \ref{fig:argon_b}.
For those setup the number of formations, their luminosity and the diameters are not influenced by the presence of the grounded ring.
The most relevant differance between those measures and the others it's that with grounded ring mean position of formations travels with a velocity of $v_A = \SI{30.49(410)}{\kilo\meter/\second}$, while when there isn't the ring we find $v_A = \SI{61.24(284)}{\kilo\meter/\second}$, an higher value.

An explanation could be that if plasma is repelled by the electrode it is attracted by a conductor at ground potential, when there isn't the grounded ring the target is the conductor at ground potential more close to the electrode, so plasma is attracted by it.

\begin{figure}
 \centering
 \subfloat[Number of formations.]{
    \includegraphics[width=0.48\textwidth]{Images/Shape/argon_de_N.png}
 }
 \hfill
 \subfloat[Mean luminosity.]{
    \includegraphics[width=0.48\textwidth]{Images/Shape/argon_de_Im.png}
 }
 
 \subfloat[Barycenter along x.]{
    \includegraphics[width=0.48\textwidth]{Images/Shape/argon_de_xb.png}
 }
 
 \subfloat[Diameter along x.]{
    \includegraphics[width=0.48\textwidth]{Images/Shape/argon_de_dx.png}
 }
 \hfill
 \subfloat[Diameter along y.]{
    \includegraphics[width=0.48\textwidth]{Images/Shape/argon_de_dy.png}
 }
 \caption{Analysis of argon plasma formations, when there is a target in front of nozzle exit, with grounded ring (set D) and without it (set E).}
 \label{fig:argon_af}
\end{figure}



%\section{Simple 1D-model}
%General desription, why and Comsol.

%\subsection{Theory}

%\subsection{Results and measure comparision}

\clearpage
\chapter{Plasma spectrum}
\label{ch:spectrometry}
%One of the fundamental characteristics of Plasma Coagulation Controller for medical applications is what species are produced and deposited during its application.
In various studies the spectrum of plasma DBD discharge in air at atmosferic pressure and ambient temperature has been observed (\cite{DBDair_Trot}, \cite{DBDAirTypicalSpec}). It presents peaks relative to reactive species from water, oxygen, nitrogen and its oxides in a wavelength range from $\num{200}$ to $\SI{880}{\nano\meter}$.

Reactive species that are thought to be involved in blood coagulation mechanisms are Reactive Oxidant Species (such as hydroxil radical \ce{OH}) and Reactive Nitrogen Species (derived from nitric oxide \ce{NO}) \cite{6153386}. This spectroscopy study is focalized on them and their precursor, i.e. on transitions relative to hydroxil, oxygen and molecular nitrogen.


\section{Optical Emission Spectroscopy}
The source produces plasma from a mixture gas of helium (or neon or argon) and air. Along free electrons, there are ions and species that collides with energetic electrons, and populate excited or metastable states with short lifetime. All reactive species partecipe in different reactions, and also in excitation and de-excitation reactions with conseguent emission of radiation. When an electron goes from state $p$ at higher energy to state $k$ of lower energy, radiation is emitted with central wavelenght $\lambda_0$. Power emitted by this radiation is given by radiant flux $d\phi_{\lambda}$ and selecting a solid angle as in figure \ref{fig:emissionfig}, is possible to define radiance $L_{\lambda}$, and intensity $I$, as in equations \ref{eq:emission}. Intensity for a radiation ultimately depends on $n(p)$, population density for state p, and Einstein Coefficient for the transition $A_{pk}$ that is typical for the transition \cite{book:291477}.
\begin{equation}
 \centering
 \begin{split}
 &\lambda_{0} = \frac{hc}{E_p - E_k} \\
 &L_{\lambda} = \frac{d^2\phi_{\lambda}}{dA \cos(\theta) d\Omega} \\
 &I = \int L_{\lambda} d\lambda = n(p) A_{pk}
 \end{split} 
 \label{eq:emission}
\end{equation}
\begin{figure}
 \centering
 \includegraphics[width = 0.45\textwidth]{Images/Spectroscopy/plasmaemission.png}
 \caption{Representation of plasma emission from a section of plasma with area $dA$, under a solid angle $d\Omega$ in a direction $\theta$.}
 \label{fig:emissionfig}
\end{figure}

Reactions that emit radiation in visible wavelenght in air are vibronic transitions where molecule goes from a vibrational state to another, with a change of vibrational quantum number $\nu$, and/or from a rotational state to another, with change of quantum number $J$ (\cite{book:137793}, \cite{wiki:vibronic}). When there is a vibrational transition, each line corresponds to different numbers $\nu'-\nu''$, these are transitions well spaced in the spectrum, easy to recognize. Rotational transitions gives birth to bands of peaks not markedly spaced, hard to resolve whitout an efficient spectrometer.

There are many reactions involving oxygen and nitrogen (see for example \cite{Kossyi_1992}), in this study only principal transition observable with our spectrometer are determined, to know dominant reactive species present in the plasma plume.

In the first section of this study plasma emission lines are recognized and studied for different discharge parameters. In the second section intensities relative to specific elements are observed for different positions along plasma plume and different gas compositions. 

\section{Line recognition}
Prototypes presented in chapter \ref{ch:electric} produce equivalent discharge conditions.
Plasma emission line recognition is done with prototype \textbf{A}, because it was the one that could be moved more easily to reach the spectrometer used in this part of the study.
A metal plate is positioned as target at a distance of $\SI{10}{\milli\meter}$ from plasma exit, as in figure \ref{fig:app1}. Helium with a flow set to $\SI{2}{\liter/\minute}$ is used to start the discharge.

The spectrometer is an IsoPlane that separates emissions with different wavelenghts with a grating. The spectrometer has a focal lenght of $\SI{320}{\milli\meter}$ and is equipped with three different gratings: $\num{150}$, $\num{1200}$ and \SI{2400}{gg/\milli\meter}, corresponding to different resolutions. %of $\num{0.26}$, $\num{0.03}$ and $\SI{0.01}{\nano\meter}$.
Light emitted by plasma is collected through a quartz lens, on the left in  figure \ref{fig:app1},%...
and travels in an optical fiber %..
connected to the spectrometer entry. At the spectrometer exit there is a CCD camera of $2048$ pixels and a count limit of $\num{65000}$.
\begin{figure}
\centering
\includegraphics[width=.6\textwidth]{Images/Spectroscopy/apparato.jpg}
\caption{Picture of the experimental setup for line recognitizion. It is possible to see the source, the metal target and the optical setup on the left. Plasma emission is collected by the lens and sent to the spectrometer.}
\label{fig:app1}
\end{figure}

Once a grating is chosen, it is possible to set the starting wavelength on the acquisition system and from there the spectrometer takes measures until the end of the CCD, resulting in a specific wavelength interval for each grating.
For each measure we select an appropriate acquisition time that allows to observe peaks with a good count number and avoid saturation.

It's important to stress out that, with this measuring method and due to complexity of plasma reactions and composition, it's not possible to extrapolate quantitative considerations between different species concentration. However it's possible to recognize the presence of certain species and make some considerations watching spectra variation with different experimental setup.


\subsection{Emission measurements}
The emissions spectrum is observed in a wavelength range from $\num{230}$ to $\SI{800}{\nano\meter}$, with standard discharge parameters: $f \SI{5}{\kilo\hertz}$ and $V_{p} = \SI{6.0}{\kilo\volt}$.
The first measurs is a rapid acquisition with the lowest resolution possible, to see intersting regions and have an idea of required exposition times. After that, the spectrometer does another acquisition with higher resolution for all wavelengths, measuring several spectra. The entire spectrum is reconstructed attaching different spectra, showed in figure \ref{fig:spectr}, where are labelled principal transitions.
For each measure a background spectrum si also taken, without plasma, to recognize peaks that are not relative to plasma emission.

The spectra are read with IDL routines \cite{GUMLEY200215} and analyzed with ROOT \emph{TSpectrum.h} library \cite{ROOT:TSpectrum}.
In each spectrum each channel is divided by the exposition time of the spectrum, evaluating the count rate at a specific wavelength.
White noise contribution is estimated as the average value from a portion of the spectrum that does not presents peaks, hence it is subtracted to count rates for each wavelength.
In the resulting spectra emission peaks are found with TSpectrum functions (where is possible to set a treshold in heigth and the general width for lines to be searched) and isolate peaks from background. The exact wavelength for each transition is found with a gaussian fit in an interval that takes into consideration the asymmetry of the peak where it's needed.

As said before, this study is focused on measure related to ROS and RNS, so in lines for \ce{NO}, \ce{OH} and \ce{N_2}.
\begin{figure}
\centering
\includegraphics[width=0.99\textwidth]{Images/Spectroscopy/spettrotot_unico_label_def.png}
\caption{Spectrum with helium flow of $\SI{2}{\liter/\minute}$, pulse parameters of $f = \SI{5}{\kilo\hertz}$ and $V_p = \SI{6.0}{\kilo\volt}$. The emission is collected at the end of the nozzle, where plasma exits in air.}
\label{fig:spectr}
\end{figure}


\paragraph{\ce{NO} lines}
In the spectrum there are two doublets for the transition $\ce{A^2\Sigma^+} \rightarrow \ce{X^2\Pi}$ with vibrational numbers (0-0) and (0-1) (\cite{Knie:166349}, \cite{VANSPRANG197955}), presented in table \ref{tab:spettroNO}. Intensities for the peaks are normalized with maximum value of $\num{1000}$ for the acquisition, the table shows as the intensities for this transition is very low. Other transitions relative to this molecule have even lower relative intensity and are not observed.
\begin{table}[h]
\centering
 \begin{tabular}{cc}
  \toprule
  $\lambda$ \text{[}\si{\nano\meter}\text{]} &\text{I [arb.u.]}\\
  \midrule
  \num{236.31(24)}  &27\\
  \num{237.00(15)}  &26\\
  \num{247.02(5)}  &28\\
  \num{247.86(12)}  &27\\
  \bottomrule
 \end{tabular}
 \caption{Peaks measured for \ce{NO}.}
 \label{tab:spettroNO}
\end{table}


\paragraph{\ce{OH} lines}
In the spectrum is observed the rotational band for transition (\ce{A^2\Sigma}, $\nu' = 0$ $\rightarrow$ \ce{X^2\Pi}, $\nu'' = 0$), $13$ lines are found \cite{doi:10.1142/S0129183100000857}.
In figure \ref{fig:OHsp} a portion of the spectrum in intersting wavelength range for \ce{OH}, in table \ref{tab:sptrOH} are presented peak values.
\begin{figure}
 \centering
 \includegraphics[width=0.6\textwidth]{Images/Spectroscopy/OH_f5t16v_2.png}
 \caption{Emission spectrum of helium discharge with $f = \SI{5}{\kilo\hertz}$ and $V_p = \SI{6.0}{\kilo\volt}$ from \num{306} to \SI{310}{\nano\meter}. It is possible to see the rotational band of \ce{OH}.}
 \label{fig:OHsp}
\end{figure}
\begin{table}
 \centering
 \begin{tabular}{cc}
  \toprule
  $\lambda$ \text{[}\si{\nano\meter}\text{]} &\text{I [arb.u.]}\\
  \midrule
  \num{306.96(1)}  &53\\
  \num{307.11(1)}  &58\\
  \num{307.29(1)}  &62\\
  \midrule                          
  \num{307.94(1)}  &142\\
  \num{308.09(1)}  &148\\
  \num{308.26(1)}  &161\\
  \num{308.43(1)}  &112\\
  \num{308.62(1)}  &46\\
  \num{308.74(1)}  &137\\
  \midrule                          
  \num{309.11(1)}  &151\\
  \num{309.22(1)}  &120\\
  \num{309.45(1)}  &36\\
  \num{309.73(1)}  &125\\
  \bottomrule
 \end{tabular}
 \caption{Peaks position and intensity measured for \ce{OH} in helium discharge with parameters $f = \SI{5}{\kilo\hertz}$, $V_p = \SI{6.0}{\kilo\volt}$.}
 \label{tab:sptrOH}
\end{table}


\paragraph{\ce{N_2} and \ce{N+_2} lines}
Measured spectrum presents several lines relative to nitrogen molecules, including the strongest, for diatomic molecule dinitrogen. The Second Positive System for \ce{N2} transition $\ce{C^3\Pi} \rightarrow \ce{B^3\Pi}$ and the First Negative System for \ce{N2+} transition $\ce{B^2\Sigma} \rightarrow \ce{X^2\Sigma}$ are observed. Their positions and peak values are presented in table \ref{tab:sptrN} (\cite{N2lab}, \cite{Britun_2007}). For \ce{N2} also a band of multiple rotational lines centered around \SI{336.58(1)}{\nano\meter} is observed.
Some of the peaks are seen in the second diffraction order, where there is more distance between lines. In figure \ref{fig:N2} are presented two portions of the spectrum showing \ce{N2} lines.
\begin{figure}
 \centering
 \subfloat[Transitions with $\Delta \nu = 2$]{
    \includegraphics[width=0.45\textwidth]{Images/Spectroscopy/N2v_f5t16_2.png}
 }
 \hfill
 \subfloat[Strongest line (0-0), $2^{\text{nd}}$ diffraction order]{
    \includegraphics[width=0.45\textwidth]{Images/Spectroscopy/N2r_f5t16_2.png}
 }
 \caption{Emission spectrum of helium discharge with $f = \SI{5}{\kilo\hertz}$ and $V_p = \SI{6.0}{\kilo\volt}$ from \num{365} to \SI{382}{\nano\meter} (a) and from \num{670} to \SI{675}{\nano\meter} (b). It is possible to see vibrational and rotational bands for \ce{N_2} transitions.}
 \label{fig:N2}
\end{figure}

\begin{table}
\centering
 \begin{tabular}{cccc}
  \toprule
                            &$\lambda$ \text{[}\si{\nano\meter}\text{]} &\text{I [arb.u.]}  &($\nu'-\nu''$)\\
  \midrule
  \multirow{3}*{\ce{N_2}}   &\num{316.03(1)}  &381  &(1-0)\\
                            &\num{337.11(1)}  &1000 &(0-0)\\
                            &\num{357.77(1)}  &722  &(0-1)\\
  \midrule
  \multirow{4}*{\ce{N_2}}   &\num{367.22(20)}  &58  &(3-5)\\
                            &\num{371.12(4)}  &172  &(2-4)\\
                            &\num{375.66(2)}  &232  &(1-3)\\
                            &\num{380.64(2)}  &423  &(0-2)\\
  \midrule
  \multirow{2}*{\ce{N_2^+}} &\num{391.50(2)}  &355  &(0-0)\\
                            &\num{427.45(2)}  &180  &(0-1)\\
  \bottomrule
 \end{tabular}
 \caption{Peaks position and intensity measured for \ce{N_2} and \ce{N+_2} in helium discharge with parameters $f = \SI{5}{\kilo\hertz}$, $V_p = \SI{6.0}{\kilo\volt}$.}
 \label{tab:sptrN}
\end{table}


\paragraph{Atomic lines}
Other lines relative to different elements are presented in table \ref{tab:sptrother}, they are \cite{NIST}:
\begin{itemize}
 \item \textbf{\ce{H_{\alpha}}} line corresponding to transition from quantum number $n=3$ to $n=2$
 \item \textbf{\ce{He}} two of the strongest lines for helium
 \item \textbf{\ce{O}} strong line of oxygen
\end{itemize}
\begin{table}
\centering
 \begin{tabular}{ccc}
  \toprule
                            &$\lambda$ \text{[}\si{\nano\meter}\text{]} &\text{I [arb.u.]}\\
  \midrule
  \ce{H_{\alpha}}           &\num{655.96(4)}  &113\\
  \midrule
  \multirow{2}*{\ce{He}}    &\num{586.94(5)}  &122\\
                            &\num{705.56(1)}  &649\\
  \midrule
  \ce{O}                    &\num{776.89(1)}  &393\\
  \bottomrule
 \end{tabular}
 \caption{Peaks position and intensity measured for \ce{H}, \ce{He} and \ce{O} in helium discharge with parameters $f = \SI{5}{\kilo\hertz}$, $V_p = \SI{6.0}{\kilo\volt}$.}
 \label{tab:sptrother}
\end{table}


\subsection{Pulse settings}
Plasma emission intensity is measured with different voltage peak values and pulse repetition rates to understand how discharge settings influence reactive species production.

As seen in chapter \ref{ch:electric} for different pulse repetition rates the electric behavior stays constant, however this parameter could still influence species production rate. Pulse repetition rate is related to the energy deposited on the ionized gas, an higher rate implies more pulses in a given time.

%Reactions that produce and recombine reactive species, and consequently density and lifetime of species, are influenced by electric field and duration of the discharge.
Emission spectrum is observed with three different parameter combinations, corresponding to different intensity of the treatment:
\begin{itemize}
 \item low: $f = \SI{5}{\kilo\hertz}$ and $\Delta t = \SI{15}{\micro\second}$
 \item medium: $f = \SI{10}{\kilo\hertz}$ and $\Delta t = \SI{10}{\micro\second}$
 \item high : $f = \SI{15}{\kilo\hertz}$ and $\Delta t = \SI{10}{\micro\second}$
\end{itemize}

For each setting plasma emission is observed along two different lines of sight:
\begin{itemize}
 \item position 1: as close as possible to the end of the nozzle, near plasma exit point from the source;
 \item position 2: close to the target, at \SI{10}{\milli\meter} from plasma exit point, where it collides with target.
\end{itemize}

Intensities for \ce{OH} and \ce{N2} species are evaluated collectively for the lines in a specific wavelength range. For \ce{OH} lines all the rotational band between $\num{306}$-\SI{309}{\nano\meter} is considered. Lines relative to \ce{N2} are separated in those between $\num{335}$-\SI{337}{\nano\meter} (rotational band and (0-0) transition) and those between $\num{368}$-\SI{382}{\nano\meter} (vibrational transitions with $\Delta \nu = 2$).

In figure \ref{fig:irel} there are measurement results for considered lines.

Intensities for \ce{OH} decreases drastically increasing the distance from the source, in position 2 values are lower than $0.1\%$ of those from position 1. \ce{OH} lines for both positions have same intensity with low and medium power setup, while intensity is lower with higher frequency. For both positions there is a similar behavior.

Also \ce{N2} intensities are related to pulse repetition rate: they decrease as the pulse rate increases. Increasing the rate from \num{5} to \SI{15}{\kilo\hertz} intensity values decrease around $60\%$ in position 1, and reach lower values for position 2.

It seems that production of both those reactive species have rates dependant on pulse repetition rates, in particular the intensity of their emission decreases at higher frequencies.
\begin{figure}
\centering
 \subfloat[\ce{OH} intensities in range $\num{306}$-\SI{309}{\nano\meter}]{
    \includegraphics[width=0.48\textwidth]{Images/Spectroscopy/If_OH.png}
 }
 \hspace{0.55\textwidth}
 \subfloat[\ce{N2} intensities range $\num{335}$-\SI{337}{\nano\meter}]{
    \includegraphics[width=0.48\textwidth]{Images/Spectroscopy/If_N2r.png}
 }
 \hfill
 \subfloat[\ce{N2} intensities range $\num{368}$-\SI{382}{\nano\meter}]{
    \includegraphics[width=0.48\textwidth]{Images/Spectroscopy/If_N2v.png}
 }
 \caption{Relative intensities of selected portions of the spectrum, for different pulse repetition rates, for position 1 and position 2.}
 \label{fig:irel}
\end{figure}


\subsection{Estimation of plasma temperatures}
From diatomic molecule's spectra it's possible to evaluate some parameters that are indicators of plasma's state: rotational temperatures for \ce{OH} and \ce{N2}, $T_{r}$, and vibrational temperature for \ce{N2}, $T_{v}$.
These parameters are estimation of the temperature at which thermal energy is comparable to the gap energy between rotational or vibrational state transitions, they can be defined as in equations \ref{eq:temperatures} where $\nu$ is the vibrational quantum number and $I$ is the quantizied moment of inertia of the molecule.
\begin{equation}
 \centering
 \begin{split}
  &T_{r} = \frac{\hslash^2}{2 k_{B} I}\\
  &T_{v} = \frac{h \nu}{k_{B}}
 \end{split}
 \label{eq:temperatures}
\end{equation}

Rotational temperatures can be considered an estimation of neutral gas kinetic temperature. Vibrational temperature gives an idea of the population of vibrational states, useful to determine chemical reactions inside plasma.

\subsubsection{Rotational temperature for \ce{OH} and for \ce{N2}}
In rotational bands the intensity of a transition for a specific wavelength is proportional to the number density popolation of upper state (equation \ref{eq:emission}), that, considering a Maxwell-Boltzmann distribution, is proportional to the temperature of the species. In equation \ref{eq:Irot} the proportionality is explicitated, with $D_{0}$ parameter that depends on number of initial molecules, partition function of the rotational state and quantum rotational numbers for upper and lower state, $S$ is the oscillator strenght specific for the molecule and $E_{r}$ depends from a constant defined by the vibrational state and from quantum rotational number for upper state \cite{MOON2003249}.
\begin{equation}
 \centering
 I = \left(\frac{2 \pi}{\lambda}\right)^4 D_0 S \exp\left(-\frac{E_r}{k_B T_r}\right)
 \label{eq:Irot}
\end{equation}

As explained before, rotational bands have many lines, not all distinguishable with the spectrometer used in this work. An approach to temperature estimation is to simulate spectra with different temperatures and to minimize differences for measured spectrum (\cite{Izarra_2000}).
In a predetermined range of temperatures spectra with different temperatures are simulated, where each line is a gaussian peak with its width that takes into consideration broadning due to thermal motion, Doppler effect and measure resolution. For each spectrum the mean square difference is evaluated and the temperature associated with the minimum difference is chosen as the one relative to the spectrum. The error is estimated taking an upper and a lower limit where the mean square difference is larger than $5\%$ of the minimum value. An example of the spectrum is shown in figure \ref{fig:Trfit}, while resulting temperatures are shown in figure \ref{fig:Trval}.
\begin{figure}
\centering
  \subfloat[\ce{OH}]{
    \includegraphics[width=0.48\textwidth]{Images/Spectroscopy/OHFit_f5t16elio.png}
  }
  \hfill
  \subfloat[\ce{N_2}]{
    \includegraphics[width=0.48\textwidth]{Images/Spectroscopy/N2rotFit_elio.png}
  }
 \caption{Example of optimal spectrum simulation for \ce{OH} and \ce{N_2} considered species.}
 \label{fig:Trfit}
\end{figure}


Estimated temperatures are compatible with each other, for each distance and for each pulse setup. It's then possible to evaluate a mean value for the species, that are $T_{r,\ce{OH}} = \SI{352(38)}{\kelvin}$ and $T_{r,\ce{N2}} = \SI{321(41)}{\kelvin}$, compatible with each other, and that, as said before, can be taken as an indicator of kinetic temperature for neutral species, so as temperature of the fluid. Those temperatures are higher then room temperature, but they are compatible with the definition of cold plasma.
\begin{figure}
\centering
  \subfloat[\ce{OH}]{
    \includegraphics[width=0.48\textwidth]{Images/Spectroscopy/TrOH2.png}
  }
  \hfill
  \subfloat[\ce{N2}]{
    \includegraphics[width=0.48\textwidth]{Images/Spectroscopy/TrN2_2.png}
  }
 \caption{Estimation of rotational temperature of \ce{OH} and \ce{N2} molecules, for different parameters setup and positions.}
 \label{fig:Trval}
\end{figure}

\subsubsection{Vibrational temperature for \ce{N_2}}
Given a set of vibrational transition lines with defined $\Delta \nu = \nu' - \nu''$, their relative intensities are correlated to each other, with a proportionality that involves vibrational temperature \cite{Britun_2007}.
With peak intensities estimated for a given transition can be made the Boltzmann graph shown in figure \ref{fig:TvN2} (a), for each experimental condition. From the Boltzmann graphs $T_v$ can be evaluated as in formula \ref{eq:Tv}.
\begin{equation}
 \begin{split}
  & log(I(\nu') = S \nu'' + I_0 \\
  & T_v [K] = \frac{10^4}{3.57 \cdot S - 0.03}
 \end{split}
\label{eq:Tv}
\end{equation}

Results are in figure \ref{fig:TvN2}. For this parameter values are compatible with each other at low and medium power pulse settings with a mean value of $T_v = \SI{3405(154)}{\kelvin}$, while a lower temperature is found for high power pulse settings $T_v = \SI{2781(322)}{\kelvin}$. It seems that whith an higher pulse repetition rate there is lower concentration of excited \ce{N_2}, produced with lower energy.
\begin{figure}
\centering
  \subfloat[Boltzmann graph for discharge with $f = \SI{5}{\kilo\hertz}$, $\Delta t = \SI{16}{\micro\second}$]{
    \includegraphics[width=0.48\textwidth]{Images/Spectroscopy/boltzman_f5t16v_2.png}
  }
  \hfill
  \subfloat[vibrational temperature of \ce{N2}]{
    \includegraphics[width=0.48\textwidth]{Images/Spectroscopy/TvN2_2.png}
  }
  \caption{(a) Boltzmann graph made from \ce{N_2} vibrational transition intensities for helium discharge with $f = \SI{5}{\kilo\hertz}$ and $V_p = \SI{6.0}{\kilo\volt}$, in position 1; (b) vibrational temperatures of \ce{N2} molecule for different pulse paramters settings and emission measured from position 1 and 2.}
  \label{fig:TvN2}
\end{figure}


\section{Line intensity analysis}
The second phase of this study is dedicated to measurements of plasma emission intensiy at different positions along the plasma plume and with different gas composition.

For this study are used a portable mini spectrometer and source prototype \textbf{B}, with an experimental setup similar to the one in chapter \ref{ch:shape}. 

\subsection{Experimental setup}
In figure \ref{fig:app2} can be seen the source and the lens focalized where plasma forms, mounted on a support with variable height. The source mounts the glass nozzle used in chapter \ref{ch:shape}, allowing to observe emission also inside it. A conductive target is placed at \SI{10}{\milli\meter} from nozzle end.
\begin{figure}
\centering
\includegraphics[width=.7\textwidth]{Images/Spectroscopy/app2_lines.png}
\caption{Picture of the experimental setup for intensity measurements. Spectrometer's lens on the left is focalized on the plasma plume, mounted on a support with variable height.}
\label{fig:app2}
\end{figure}

It is used a mini spectrometer \emph{Hamamatsu C10082CAH}, with resolution of \SI{1}{\nano\meter}, spectral range from \SI{150}{\nano\meter} to \SI{850}{\nano\meter} and a CCD sensor \emph{S10420-1106} with \num{2048} pixels. The efficiency of the spectrometer for different wavelengths is evaluated with measurements of emission from a lamp with known emission spectrum. With this calibration is possible to convert the number of counts in the number of photons emitted by plasma for each wavelength.

Helium, neon or argon are used to produce plasma, spectra from those three gasses is shown in figure \ref{fig:specs}. Given the low resolution of this spectrometer, it is not possible to identify all the peaks that are observed. Each spectrum presents the same emission lines and, in addition, lines relative to the element used to start the discharge. It is possible to identify four wavelength ranges relative to four elements:
\begin{itemize}
 \item \ce{N_2} lines from \num{320} to \SI{480}{\nano\meter}: one singlet and several doublets relative to the SPS described before in this chapter; 
 \item one \ce{He} line at \SI{706.6}{\nano\meter}, observed only for helium plasma;
 \item \ce{Ne} lines from \num{410} to \SI{780}{\nano\meter}, several lines observed only for neon plasma with intensity generally higher than \ce{N_2} lines;
 \item \ce{Ar} lines from \num{660} to \SI{850}{\nano\meter}, several lines observed only for argon plasma.
\end{itemize}
\begin{figure}[h]
 \centering
 \subfloat[Spectrum with helium gas.]{
    \includegraphics[width=0.65\textwidth]{Images/Spectroscopy/B_elio_sp.png}
  }
  
  \subfloat[Spectrum with neon gas.]{
    \includegraphics[width=0.65\textwidth]{Images/Spectroscopy/B_neon_sp.png}
  }
  
  \subfloat[Spectrum with argon gas.]{
    \includegraphics[width=0.65\textwidth]{Images/Spectroscopy/B_argon_sp.png}
  }
  \caption{Plasma emission spectrum measured at the end of the nozzle, with different gasses used to start the discharge.}
  \label{fig:specs}
\end{figure}


Emission for different elements is evaluated selecting, for each element, some of the higher peaks in all spectra between those observed:
\begin{itemize}
 %\item UVB lines : doublet centered at \SI{293.2}{\nano\meter} and \SI{295.5}{\nano\meter} that could be emitted by \ce{NO} molecules, but we can not be sure due to spectrometer's low resolution;
 \item \ce{N_2} : three lines centered at \SI{353.6}{\nano\meter} (transition $(1-2)$), \SI{380.5}{\nano\meter} (transition $(0-2)$) and \SI{405.9}{\nano\meter} (transition $(0-3)$);
 \item \ce{He} : one line at \SI{706.6}{\nano\meter};
 \item \ce{Ne} : one doublet at \SI{638.3}{\nano\meter} and \SI{640.6}{\nano\meter}, a single line at \SI{703.5}{\nano\meter} and a single line at \SI{725.5}{\nano\meter};
 \item \ce{Ar} : three different lines at \SI{696.5}{\nano\meter}, \SI{763.5}{\nano\meter} and \SI{811.5}{\nano\meter}.
\end{itemize}


\subsection{Position}
The experimental setup allows to measure plasma emission at different heights along the nozzle axis, emission intensity evaluated in four different positions:
\begin{itemize}
 \item \SI{-5}{\milli\meter}, inside the nozzle
 \item \SI{0}{\milli\meter}, at the end of the nozzle
 \item \SI{5}{\milli\meter}, between nozzle and target
 \item \SI{10}{\milli\meter}, right before target position
\end{itemize}

Results are in figure \ref{fig:Irel_pos}. Total intensities are normalized to compare emissivity variation for different positions, while relative intensities are divided for the respective total intensity.
With helium emission is lower inside the nozzle, increase at the exit and decreases slightly moving towards the target. For other gasses emission decreases increasing distance from the source.
For every gas, emission from \ce{N_2} increases until it reaches a constant value.
Lines from specific gasses decreases distant from the nozzle, with slope different for different gasses.
It's intersting to note how neon emission decreases linearly, while \ce{N_2} emission increases and it's an effect that is observed visually on the plume: going from the nozzle to the target, plasma colour goes to red (neon emission) to violet (nitrogen emission).
\begin{figure}
\centering
  \subfloat[Total emission intensity.]{
    \includegraphics[width=0.48\textwidth]{Images/Spectroscopy/Itot_pos2.png}
  }
  
  %\subfloat[UVB lines]{
  %  \includegraphics[width=0.48\textwidth]{Images/Spectroscopy/Irel_NO_pos.png}
  %}
  
  \subfloat[\ce{N_2} lines]{
    \includegraphics[width=0.48\textwidth]{Images/Spectroscopy/Irel_N2_pos2.png}
  }
  \hfill
  \subfloat[\ce{He} lines, \ce{Ne} lines, \ce{Ar} lines.]{
    \includegraphics[width=0.48\textwidth]{Images/Spectroscopy/Irel_gas_pos2.png}
  }
\caption{Axial profile of total intensities (a) and relative intensities for selected portions of the spectrum (b-c), with different starting gas. At \SI{0}{\milli\meter} the lens points at the end of the nozzle, at \SI{10}{\milli\meter} at a metal target. Relative intensities in (c) are for lines corresponding to the element that used as starting gas. Relative intensities take into consideration spectrometer's efficiency and total emission for each position.}
 \label{fig:Irel_pos}
\end{figure}

It can be concluded that inside the nozzle the emission depends from the gas used to start the plasma, going outside the emission due to nitrogen species becames dominant in the discharge.


\subsection{Gas composition}
Relative intensity for selected lines is analyzed for different gas composition. Helium is mixed with argon or nitrogen using specific flowmeters, that allow to add up to \SI{0.2}{\liter/\minute}, with resolution of \SI{0.01}{\liter/\minute}.
With an helium flow to \SI{2}{\liter/\minute} is possible to have gas mixtures where argon and nitrogen have a maximum percentage of $\num{10}\%$.

The spectrometer is pointed at the end of the nozzle and collects spectra for three different gas concentrations, measuring intensities shown in figure \ref{fig:Irel_flow}.

When nitrogen is added to the gas, total emission intensity lowers, while nitrogen emission does not change. It is possible that the increase in nitrogen concentration is too low if compared to the quantity naturally present in air and consequently there is not a significative increase of nitrogen emission at the end of the nozzle.
In this position the gas is a mix between the gas used to start the discharge and air, where nitrogen is approximately the $70\%$ of air. The percentage of gas in air is a function of the gas flow, for example it is possible to hypotize that at the end of the nozzle the percentage of discharge gas and air is equal. With a gas flow of pure helium, nitrogen would be the $35\%$ of the total gas density. If the helium flow is substituted with $10\%$ of nitrogen and $90\%$ helium, as in the experiment condition, nitrogen would be the $40\%$ of total gas density. If the relative emission of nitrogen is proportional to the fraction of nitrogen in the gas, the increase in the emission would be a value around $5\%$, an increment that could be not resolvable with the apparatus used in this experiment. Further measurements of nitrogen emission could verify this hypothesis and, in that case, they could be used to estimate nitrogen percentage for different positions inside the nozzle. 


When argon is added to helium, total emission intensity increases, while relative emission from elements other than argon decreases slightly. It means that the only variation when there is a percentage of argon is that the emission relative to this element adds t othe emission of helium and the emission relative to other elements stays unchanged.
\begin{figure}
\centering
  \subfloat[Total emission intensity.]{
    \includegraphics[width=0.48\textwidth]{Images/Spectroscopy/Itot_flux2.png}
  }
  
  %\subfloat[UVB lines]{
  %  \includegraphics[width=0.48\textwidth]{Images/Spectroscopy/Irel_NO_flux.png}
  %}
  
  \subfloat[\ce{N_2} lines]{
    \includegraphics[width=0.48\textwidth]{Images/Spectroscopy/Irel_N2_flux2.png}
  }
  \hfill
  \subfloat[\ce{He} line]{
    \includegraphics[width=0.48\textwidth]{Images/Spectroscopy/Irel_He_flux2.png}
  }
  
  \subfloat[\ce{Ar} line for argon gas mixture]{
    \includegraphics[width=0.48\textwidth]{Images/Spectroscopy/Irel_Ar_flux2.png}
  }
\caption{Behaviour of total intensities (a) and relative intensities for selected portions of the spectrum (b-c-d), changing the composition of the gas. A flow of \SI{0.2}{\liter/\minute} corresponds to the $\num{10}\%$ of the total gas flow. Relative intensities take into consideration spectrometer's efficiency and total emission for each gas mixture.}
 \label{fig:Irel_flow}
\end{figure}


\begin{comment}
La determinazione delle specie reattive prodotte dal plasma a pressione atmosferica può essere effettuata tramite misure spettrometriche della sorgente in funzione su un bersaglio.
Gli effetti del trattamento al plasma si pensano dovuti alla presenza di ROS e RNS, quindi vengono raccolte misure nel range di lunghezze d'onda utili ad osservare le emissioni di molecole di \ce{OH} (\SI{305.00}-\SI{313.00}{\nano\meter}), \ce{N_{2}} (\SI{280.00}-\SI{500.00}{\nano\meter}) e \ce{NO} (\SI{220.0}-\SI{290.0}{\nano\meter}) (vedi articoli).

I prototipi di sorgente sviluppati permettono di variare il range dei parametri di funzionamento, in modo da modulare l'intensità del trattamento. Per verificare come cambia lo spettro in base alle diverse modalità di funzionamento, si osservano la variazione nell'intensità delle emissioni al cambiare di frequenza e tempo di apertura del circuito.

Si vuole osservare l'intensità relativa delle righe al variare del gas in ingresso, quindi la sorgente viene azionata con diverse miscele di gas. La misura standard viene effettuata con flusso di \ce{He}, nella solita modalità di funzionamento della sorgente. Vengono poi predisposte due ulteriori modalità di misura dove il gas viene fatto gorgogliare in una soluzione di acqua o di ammoniaca prima dell'inserimento all'uscita della sorgente, per arricchire i prodotti delle reazioni di ioni contenenti, rispettivamente, ossigeno o azoto.

A partire dalla forma delle righe di alcune specie molecolari, è inoltre possibile stimare la temperatura rotazionale alla quale avviene l'emissione misurata. Le emissioni dovute alla molecola \ce{OH} o \ce{N_2} sono varie righe dalle intensità variabili a seconda della temperatura delle molecole, misurando l'intensità relativa dei picchi si può ricavare la temperatura rotazionale delle molecole.

\section{Setup di acquisizione}

Per la sorgente vengono utilizzati sia il prototipo precedente, \textbf{prototipo 1} , sia il prototipo svilupato durante questo lavoro, \textbf{prototipo 2}.

Vengono inoltre provate tre diverse modalità di funzionamento della sorgente, variando frequenza e tempo di chiusura del circuito, mostrate in Tabella \ref{tab:setsorgente}.
Il flusso del gas di elio viene mantenuto pari a \SI{2}{\liter/\minute}.

\begin{table}
 \centering
 \begin{tabular}{cc}
 \toprule
 $f$ [\si{\kilo\hertz}]  &$\Delta t$ [\si{\micro\second}]\\
 \midrule
 5  &15\\
 10 &10\\
 15 &10\\
 \bottomrule
 \end{tabular}
 \caption{Parametri di funzionamento utilizzati per le diverse misure.}
 \label{tab:setsorgente}
\end{table}



\section{Presentazione ed analisi misure}
Per entrambe le sorgenti, l'analisi prevede il riconoscimento delle emissioni misurate, il confronto delle intensità variando distanza e parametri di funzionamento della sorgente, il confronto delle intensità variando il gas immesso, la stima delle temperature rotazionali delle molecole \ce{OH} e \ce{N_2}.

\subsection{Riconoscimento righe}
L'output dello spettrometro viene letto tramite routine IDL ed il riconoscimento dei picchi viene effettuato tramite la classe TSpectrum presente nelle librerie ROOT. Ogni misura viene confrontata con uno spettro di background preso con lo stesso tempo di acquisizione e sorgente spenta, riuscendo così ad escludere i picchi dovuti al fondo.
In figura \ref{fig:spettrotot} si vedono le principali righe estrapolate dalle acquisizioni, in particolare si vedono le righe relative ad \ce{NO}, \ce{OH}, \ce{H_{2}} e \ce{N_{2}}, tabulate in Tabella \ref{tab:spettrotot} .

\begin{figure}
\centering
\includegraphics[width=0.99\textwidth]{Immagini/spettrotot_unico_label_def.png}
\caption{Spettro acquisito con condizioni di misura standard ($f = \SI{5}{\kilo\hertz}$ e $\Delta t = \SI{16}{\micro\second}$), obiettivo puntato vicino l'uscita del gas dalla sorgente.}
\label{fig:app}
\end{figure}


L'acquisizione migliore, nella quale vengono riconosciuti più picchi, è quella mostrata in Figura \ref{fig:spettrotot}, corrispondente alla posizione 1 e condizioni standard di misura, $f = \SI{5}{\kilo\hertz}$.

\paragraph{}
Per verificare gli effetti di una diversa frequenza  nel funzionamento della sorgente vengono confrontate le intensità relative alle righe di \ce{OH} e \ce{N_2}, sommando i conteggi per le varie porzioni di spettro. Non vengono prese in considerazione le righe del gruppo relativo all'\ce{NO} in quanto troppo deboli. In Tabella \ref{tab:irel_1} i risultati, dove vengono confrontate i conteggi alle varie frequenze rispetto i conteggi ottenuti nelle condizioni standard di lavoro, $f = \SI{5}{\kilo\hertz}$. Si nota un calo evidente nei conteggi, crescente con l'aumentare della frequenza di lavoro.

Allo stesso modo viene variata la posizione dell'obbiettivo, dalla posizione 1, puntato all'uscita della sorgente, alla posizione 2, puntato all'uscita del bersaglio. I risultati sempre in Tabella \ref{tab:irel_1}, dove vengono presentati i conteggi nella posizione 2 rispetto i conteggi nella posizione 1. Viene trovato un effetto diverso sulle diverse specie, le emissioni di \ce{OH} diminuiscono molto, mentre quelle relative l'\ce{N_2} diminuiscono in maniera inferiore.

\begin{table}
 \centering
 \begin{tabular}{lcc}
 \toprule
                            &\ce{OH}  &\ce{N_2}\\
 \midrule
 f = \SI{5}{\kilo\hertz}    &1.00     &1.00\\
 f = \SI{10}{\kilo\hertz}    &0.95    &0.81\\
 f = \SI{15}{\kilo\hertz}    &0.62    &0.63\\
 \midrule
 posizione 1                &1.00        &1.00\\
 posizione 2                &0.10     &0.82\\
 \bottomrule
 \end{tabular}
 \caption{Intensità relative delle porzioni interessanti di spettro, al variare di frequenza e posizione, per il prototipo 1.}
 \label{tab:irel_1}
\end{table}


\subsection{\ce{He} + \ce{H_2O} e \ce{He} + \ce{NH_3}}
Vengono misurati gli spettri di emissione in condizioni di lavoro standard, $f = \SI{5}{\kilo\hertz}$, posizione 1, trattando il flusso di elio prima di inserirlo nella sorgente.
Per il prototipo 1, l'arricchimento di specie reattive dell'ossigeno o dell'azoto avviene facendo gorgogliare l'elio in una soluzione di acqua o ammoniaca, mantenendo un flusso in uscita dalla bombola sempre di \SI{2}{\liter/\minute}.
Le misure così acquisite presentano un'intensità molto minore per tutte le righe dello spettro e non vi sono variazioni rilevanti per le lunghezze d'onda interessate, cioè le emissioni relative ad \ce{OH} per la soluzione di acqua e relative ad \ce{N_2} per la soluzione di ammoniaca.

\subsection{Stima temperatura}
Lo spettro di emissione della molecola \ce{OH} è della forma presentata in Equazione \ref{eq:fitOH} (vedi articolo):
\begin{equation}
\centering
\begin{split}
&I_i (T) = I_{0,i} \exp{-\frac{E_n (T-T_0)}{T_0 T}} \\
&S_i(\lambda,T) = \frac{I_i(T)}{\sigma \sqrt{2\pi}} \exp{\frac{(\lambda - \lambda_i)^2}{2\sigma^2}}\\
&S(\lambda,T) = \sum_i S_i(\lambda,T)
\end{split}
\label{eq:fitOH}
\end{equation}

dove $I_i$ è l'intensità ad una determinata energia $E_n$ e $I_{0,i}$ è l'intensità misurata ad una temperatura di riferimento $T_0$. Le $S_i$ sono la convoluzione di $I_i$ con una distribuzione gaussiana nelle $\lambda$ e lo spettro finale $S$ sarà la somma degli spettri relativi la singola transizione.
Simulando diversi spettri nelle varie temperature, sarà possibile stimare la temperatura rotazionale degli ioni presenti nel gas.

La procedura consiste nel simulare diversi spettri $S(\lambda,T)$ come presentati in \ref{eq:fitOH}, calcolare gli scarti quadratici medi rispetto le misure e ricavare la temperatura dalla simulazione migliore. Tipicamente sono stati simulati spettri con $200$ temperature diverse, su un range di temperature possibili variabili, a seconda della misura considerata.
L'errore sulla stima della temperatura viene calcolato come media della differenza tra la temperatura minima che avesse uno scarto quadratico medio fino al $5\%$ superiore rispetto al minimo e la differenza della temperatura massima con le stesse condizioni.

In figura \ref{fig:fitT} sono presentati due esempi di fit delle emissioni interessate. In \ref{tab:Trot} vengono presentate le diverse temperature ottenute, compatibili tra loro.

\begin{figure}
 \centering
 \subfloat[][\ce{OH}]
  {\includegraphics[width=.8\textwidth]{Immagini/OHFit_f5t16elio.png}}
 \\
 \subfloat[][\ce{N_2}]
  {\includegraphics[width=.75\textwidth]{Immagini/N2rotFit_elio.png}}
 \caption{Fit delle emissioni con spettri simulati, misure prese con il prototipo 1, in condizioni standard, posizione 1.}
 \label{fig:fitT}
\end{figure}

\begin{table}
 \centering
 \begin{tabular}{cc}
 \toprule
            &T [\si{\kelvin}]\\
 \midrule
  \ce{OH}   &$336 \pm 30$\\
  \ce{N_2}  &$322 \pm 41$\\
 \bottomrule
 \end{tabular}
 \caption{Stima delle temperature di rotazione delle molecole.}
 \label{tab:Trot}
\end{table}
\end{comment}

\begin{comment}
La misura viene effettuata tramite uno spettrometro IsoPlane dalla lunghezza focale di \SI{320}{\milli\meter}, con tre diversi reticoli: \SI{150}, \SI{1200} e \SI{2400}{g/\milli\meter}, corrispondenti alle risoluzioni di ... .
La risoluzione maggiore viene utilizzata per acquisire le righe \ce{OH} e \ce{N_{2,\text{rot}}}, mentre per acquisire lo spettro totale vengono utilizzati i reticoli a piccola e media risoluzione.

Lo spettrometro è accoppiato ad una telecamera PIXIS di $2048 \times $ ... pixels quadrati dal lato di ... \si{\micro\meter}, con un massimo di \SI{65000} conteggi per il singolo canale.
La luce viene raccolta da una lente in quarzo di focale ... e diametro ... , portata da una fibra ottica dallo spessore di ... e lunghezza di ... e collegata all'entrata dello spettrometro.

Per l'acquisizione viene posizionata la sorgente in funzione a distanza di \SI{1}{\centi\metre} dal bersaglio in metallo (collegato a terra) con l'ottica focalizzata sul flusso di gas, come in foto \ref{fig:app}. Vengono distinte due posizioni dell'ottica:
\begin{itemize}
 \item posizione 1 = obiettivo sull'uscita della sorgente
 \item posizione 2 = obiettivo sul punto di impatto del plasma sulla sorgente, ad \SI{1}{\centi\meter} dall'uscita della sorgente
\end{itemize}
.

\begin{table}
\centering
 \begin{tabular}{ccc}
  \toprule
                            &$\lambda$ \text{[}\si{\nano\meter}\text{]} &\text{I [arb.u.]}\\
  \midrule
  \multirow{4}*{\ce{NO}}    &\num{236.31(24)}  &27\\
                            &\num{237.00(15)}  &26\\
                            &\num{247.02(5)}  &28\\
                            &\num{247.86(12)}  &27\\
  \midrule
  \multirow{2}*{\ce{OH}}    &\num{308.3(1)}  &106\\
                            &\num{309.1(1)}  &113\\
  \midrule
  \multirow{5}*{\ce{N_2}}   &\num{316.03(1)}  &381\\
                            &\num{337.11(1)}  &1000\\
                            &\num{357.77(1)}  &722\\
                            &\num{367.22(20)}  &58\\
                            &\num{371.12(4)}  &172\\
                            &\num{375.66(2)}  &232\\
                            &\num{380.64(2)}  &423\\
  \midrule
  \multirow{2}*{\ce{N_2^+}} &\num{391.50(2)}  &355\\
                            &\num{427.45(2)}  &180\\
  \midrule
  \ce{H_{\alpha}}           &\num{655.96(4)}  &113\\
  \midrule
  \multirow{2}*{\ce{He}}    &\num{586.94(5)}  &122\\
                            &\num{705.56(1)}  &649\\
  \midrule
  \ce{O}                    &\num{776.39(1)}  &393\\
  \bottomrule
 \end{tabular}
 \caption{Picchi rilevanti nello spettro di emissione del prototipo 1, condizioni di misura standard, posizione 1.}
 \label{tab:spettrotot}
\end{table}


Per ogni misura viene stabilito un tempo di acquisizione idoneo ad avere un numero ottimale di eventi, evitando la saturazione dei singoli canali. Una volta stabilito il tempo viene eseguita una misura di fondo, a sorgente spenta, e successivamente viene avviata la misura con sorgente attiva.
\end{comment}

\clearpage
\chapter{Plasma power estimate}
\label{ch:temperature}
Direct application of plasma leads to deposition of accelerated particles on target. For non thermal blood coagulation it's fundamental that on biological tissues that a temperature increase due to this application must be below dangerous limits.

With the supposition that the main heating mechanism of our source is by convection, it is possible to estimate heat transfer rate, i.e. power, due to plasma application on a target. During this work only inorganic targets are utilized, while PCC will be applied to samples of blood and to living subjects, where heating effects are different if compared to inorganic matter. Heat transfer in living biological tissue is a complicated combination of thermal conduction, convection, perfusion of blood and metabolic heat generation \cite{ShitTemp}. The hypotesis behind the estimate of plasma power in this work is that the interaction between heat transfer mechanisms in biological tissues and heat deposition mechanism of our source is negligible. Under those conditions, with this power estimate will be possible to evaluate the temperature increase due to plasma application on every target once their thermal characteristics are know (for a review of thermal conduction parameters in biological tissues see \cite{biotissues}).
Even if the studied case is particular and has many limits, this analysis gives a rough estimate of plasma power dependencies from application parameters such as pulse repetition rate or distance between source and target.

\section{Experimental setup}
Plasma heat power is estimated measuring how much temperature increases with the application of plasma on a target with known heat capacity. The use of a thermocamera allows to study target's temperature profile in its entirety, visualizing also heat conduction on target's borders.
An object at thermal equilibrium in a temperature range around \SI{300}{\kelvin} emits radiation in the long-wavelength infrared region (from \num{8} to \SI{15}{\micro\meter}). This radiation can be collected and it's intensity measured by a bolometer, evaluating the temperature of the object from it's infrared emission, as in common thermal cameras (\cite{Gade2014}). The detector more used is an uncooled microbolometer VOx with an array of pixels as in figure \ref{fig:microbol}. Incident radiation strikes a material that has an absorption peak in infrared wavelenght, temperature increases changing the electrical resistance of the circuit and the resulting current intensity is measured and associated to collected radiation intensity.
\begin{figure}
 \centering
 \includegraphics[width=0.4\textwidth]{Images/Temperature/Microbolometer.png}
 \caption{Rapresentation of pixel in a microbolometer detector. When incident radiation arrives on the pixel, temperature increases and electric resistance of the circuit changes, giving rise to a current intensity proportional to radiation intensity.}
 \label{fig:microbol}
\end{figure}


\paragraph{Camera}
The target is observed by a termographic camera \emph{FLIR A655sc} (\cite{flircamera}) with a spectral range of $\num{7.5}-\SI{14}{\micro\meter}$, resolution $\num{640}\times\num{480}$ and detector pitch \SI{17}{\micro\meter}, equipped with a lens with focal \SI{41.3}{\milli\meter} (field of view \ang{15}).
Temperature evolution is a phenomenon with charateristic time of several seconds, to have a reasonable resolution frame rate acquisition is set at \num{2} FPS.

The temperature conversion between intensity of the radiation and temperature is done by \emph{FLIR} software, setting the appropriate emissivity.

\paragraph{Source}
The source is the latest one, prototype \textbf{B} in chapter \ref{ch:electric}, at different pulse repetition frequencies $f$ and voltage peak values $V_p$. Plasma formation and deposition it's not a continuos phenomenon, it happens in correspondance of voltage pulses, with a charateristic time around \SI{1}{\micro\second}, see chapter \ref{ch:shape}. Given a time interval, $f$ defines the number of pulses that arrives in that interval.

The gas used to produce plasma is helium, with flow of \SI{2}{\liter/\minute}.

\paragraph{Target}
The target is an aluminium disk, with radius \SI{7}{\milli\meter} and height \SI{1}{\milli\meter}, on a plastic support with width \SI{5}{\milli\meter}.
Aluminium has specific heat capacity $c_p = \SI{897}{\joule/\kilogram\kelvin}$ and density $\rho = \SI{2700}{\kilogram/\meter^3}$, corresponding to target heat capacity $C = \SI{0.373(5)}{\joule/\kelvin}$.

Target is positioned at a distance of \SI{210}{\milli\meter} from camera lens. At those conditions, once the image is on focus, pixel to pixel distance on acquired frames is \SI{86.4}{\micro\meter}.

To verify relation between power and application distance the target is positioned at three distance values from the en of the source head: \num{5.5}, \num{7.5} or \SI{9.5}{\milli\meter}.

\paragraph{Measurements procedure}
Ultimately measurements are done with different voltage pulse repetition frequencies $f$, voltage peak values $V_p$ and distances between target and source $d$.
Once those parameters are set, the measure follows an approximative timeline:
\begin{itemize}
 \item acquisition of at least \num{4} frames for background evaluation in \SI{2}{\second}
 \item start of gas flow for a minimum of \SI{5}{\second}
 \item start of the discharge for a minimum of \SI{30}{\second}
 \item stop of the discharge and acquisition with only gas flow for \SI{5}{\second}
 \item stop of gas flow
\end{itemize}

Start and end time for those phases can be observed directly from measures.

\section{Temperature profile}
Every frame gives information on the temperature of every pixel of acquired image at a given time, as in figure \ref{fig:exframe}. To find the interested temperature difference, background temperature before discharge is evaluated as the average value acquired before gas opening and subtracted to pixel value in other frames. In a given direction the temperature profile is shown in figure \ref{fig:exframe}.
\begin{figure}
 \centering
 \subfloat[Frame $t = \SI{25}{\second}$.]{
    \includegraphics[width=0.48\textwidth]{Images/Temperature/f5t4d4_esframe_25s_2.png}
 }
 \hfill
 \subfloat[$\Delta T$ profile along x direction at $t = \SI{25}{\second}$]{
    \includegraphics[width=0.48\textwidth]{Images/Temperature/f5t4d4_esprofilo_25s_2.png}
 }
 \caption{Example of measurement during discharge with parameters $f = \SI{5}{\kilo\hertz}$, $V_p = \SI{7.2}{\kilo\volt}$ and $d = \SI{5.5}{\milli\meter}$.}
 \label{fig:exframe}
\end{figure}


Source is pointed to target center, for frames acquired with the discharge can be seen clearly where the temperature rises and can be estimated the application position. A square window with edge \SI{15}{\milli\meter} si centered around the point with maximum temperature, to see the temperature increase on all the target. A gaussian interpolation inside the square allows to find the center position, compatible in each measurements set, at coordinates $x_{c} = \SI{25.06(20)}{\milli\meter}$, $y_{c} = \SI{21.17(20)}{\milli\meter}$.

It's interesting to note that when there is only gas flowing on the target, temperature increases radially from the center, in an area larger then the target area, while when there is plasma, temperature increases rapidly in an area compatible with the target dimensions. It's possible that gas before ionization creates a gas layer with a radius larger then the target radius, but considering temperature increase only around the target it doesn't gives problems in heat power evaluation due to discharge.

From each of those frames the average temperature difference and it's evolution are evaluated as shown in figure \ref{fig:Tavg}. This graph allows to extrapolate time intervals for gas flow and discharge: where temperature changes rapidly there is gas opening, source activation, discharge end and gas flow stop. In figure \ref{fig:dis_profiles} an example of how temperature profile changes during discharge.
\begin{figure}
 \centering
 \includegraphics[width=0.6\textwidth]{Images/Temperature/f5t4d4_Tavg_lines.png}
 \caption{Average increase in temperature for every frame at different times, on a square around the target. There are three phases highlighted: gas phase before source activation, plasma phase and gas phase after source power off.}
 \label{fig:Tavg}
\end{figure}

\begin{figure}
 \centering
 \includegraphics[width=0.6\textwidth]{Images/Temperature/f5t4d4_profilo_tempi_zoom.png}
 \caption{Evolution of temperature profile in x direction, with a zoom on the region around the target, for different times and phases: $t = \SI{2}{\second}$ before gas opening, $t = \SI{6}{\second}$ after gas opening and before discharge, $t = $\num{10}, \num{20}, \SI{30}{\second} during discharge. }
 \label{fig:dis_profiles}
\end{figure}


\section{Power estimation}
During the plasma phase, temperature increase in target can be defined as the sum of power deposition plus thermalization contributions (see \cite{Pimazzoni2018}), as in equation \ref{eq:dT}, where $m$ is the mass of the portion of target considered, $c_p$ is aluminium specific heat capacity, $P$ is power deposited by plasma on target, $T_{0}$ is a temperature limit and $k$ a costant parameter dependant from the target.
\begin{equation}
 m c_p \frac{dT}{dt} = P - k(T-T_{\text{lim}})
 \label{eq:dT}
\end{equation}

The anlytical result of this equation is an exponential function, as in figure \ref{fig:dTfit} for the center pixel in the image. The downside of the exponential fit it's that isn't possible to evaluate the contribution of power deposition from it. From temperature variation for temperatures close to $T_{0}$, where the thermalization contribute is neglectable, is possible to extrapolate parameter $P$ with a linear fit. The power parameter resulting from this interpolation is not a precise value because there are also thermalization contributes, but it is useful to study its behavior varying other parameters.
\begin{figure}
 \centering
 \includegraphics[width=0.6\textwidth]{Images/Temperature/f5t4d4_fits.png}
 \caption{Average increase of temperature for a restricted time ranges inside plasma phase, with exponential and linear interpolations. The exponential fit shows the analytical solution for equation \ref{eq:dT}, the linear fit is the one from wich power is extrapolated.}
 \label{fig:dTfit}
\end{figure}

To assure that temperature variation are observed only inside the target, only pixels within a radius of \SI{7}{\milli\meter} from the center of plasma application are considered. The maverage temperature variation in those conditions follows equation \ref{eq:dTpix}, where $\rho$ is aluminium density, $h$ is target width, $A_p$ is target area, $T_{\text{avg}}$ is the average temperature of a pixel, $T_p$ is the temperature of a pixel.
\begin{equation}
 m c_{p} T_{\text{avg}}(t) = \rho h A_{p}\frac{\sum_{\text{pixels}} T_{p}}{N_{\text{pixels}}} = P t + T_{0}
 \label{eq:dTpix}
\end{equation}


Plasma deposition isn't continuos in time but it happens within time intervals, as described in chapter \ref{ch:shape}.
Power evaluated considering the total time interval of plasma application is an effective power $P_{\text{eff}}$ that doesn't take into consideration the real time interval of interaction between plasma and target, a fraction of the total one. An estimate of the power associated with a single pulse $P_{\text{pulse}}$ can be found considering the time interval for the interaction of every pulse $\Delta t_{\text{pulse}}$ multiplied by the rate of pulses during the discharge $f$, as in equation \ref{eq:power}. An appropriate time pulse duration is $\Delta t_{\text{pulse}} = \SI{1}{\micro\second}$, as shown in chapter \ref{ch:shape}.
\begin{equation}
 P_{\text{pulse}} = \frac{P_{\text{eff}}}{\Delta t_{\text{pulse}} \, f}
 \label{eq:power}
\end{equation}


\paragraph{Pulse rate dependency}
Power parameter is evaluated with different pulse repetition rate, voltage peak value set at $V_{p} = \SI{7.2}{\kilo\volt}$ and distance between source and target $d = \SI{5.5}{\milli\meter}$. In figure \ref{fig:Pfr} there are effective power during the discharge and the pulse power relative to a single pulse. As expected the effective power is proportional to frequency of pulses: greater frequency means more pulses in a given time. It's instead intersting to note that power for a single pulse decreases with increasing pulse rate. 
An explanation could be that for different pulse repetition rate the time interval relative to the single pulse, i.e. the time interval during which we see plasma bullets, decreases from $\Delta t_{\text{pulse}} = \SI{1}{\micro\second}$ for $f = \SI{5}{\kilo\hertz}$ to values around $\Delta t_{\text{pulse}} = \SI{500}{\nano\second}$ for $f = \SI{25}{\kilo\hertz}$. This result could imply that for higher frequencies bullets have increasing speed.
\begin{figure}
 \centering
 \subfloat[Average power during discharge.]{
    \includegraphics[width=0.48\textwidth]{Images/Temperature/Peff_fr_2.png}
 }
 \subfloat[Power relative to a single voltage pulse.]{
    \includegraphics[width=0.48\textwidth]{Images/Temperature/Ppulse_fr_2.png}
 }
 \caption{Plasma power deposition on target for different pulse repetition rates.}
 \label{fig:Pfr}
\end{figure}


\paragraph{Voltage peak value dependency}
Temperature increase is observed with distance between source and target $d = \SI{5.5}{\milli\meter}$ and pulse repetition rate $f = \SI{5}{\kilo\hertz}$, for different voltage values. From figure \ref{fig:PVp} it is possible to see that power for the single pulse is correlated with the voltage peak value and increases rapidly when voltage increases. For voltage values from \SI{5.0}{\kilo\volt} to \SI{9.2}{\kilo\volt} temperature increase rate goes from \SI{7.8}{\milli\kelvin/\second} to \SI{85.2}{\milli\kelvin/\second}. With the higher voltage it is observed a temperature increase up to \SI{2.2}{\kelvin} in \SI{25}{\second} of plasma application. This behavior is expected because with an higher peak voltage value there is more electrical power transfered from the source and bullets have higher velocity, as described in chapter \ref{ch:shape}.
\begin{figure}
 \centering
 \includegraphics[width=0.6\textwidth]{Images/Temperature/Ppulse_Vp_2.png}
 \caption{Plasma power of a single pulse for different voltage peak values.}
 \label{fig:PVp}
\end{figure}

\paragraph{Distance dependency}
The influence of the distance between target and source is measured setting voltage peak at $V_{p} = \SI{7.2}{\kilo\volt}$ and $f = \SI{5}{\kilo\hertz}$. In figure \ref{fig:Pd} it is possible to see how power decreases with incresing distance, as expected because for higher distances between source and target, bullets lose more energy before the impact and will transfer less energy to the target. For distance going from \SI{5.5}{\milli\meter} to \SI{7.5}{\milli\meter} the pulse power loss rate is of \SI{0.28(5)}{\watt/\milli\meter}.
\begin{figure}
 \centering
 \includegraphics[width=0.6\textwidth]{Images/Temperature/Ppulse_d_2.png}
 \caption{Plasma power of a single pulse for different target distances.}
 \label{fig:Pd}
\end{figure}

\clearpage
\chapter{Conclusions}
\label{ch:conclusions}
In this work a source for Cold Atmospheric Plasma has been developed, charachterized and utilized to study formation and propagation of plasma produced with a DBD reactor.
The source presents an electrode covered in dielectric placed inside a nozzle where there is neutral gas flowing. The electrode applies an electric field to the neutral gas producing plasma that is expelled in air.

The first phase of this work was to develop the components for a new prototype of the source and assemble them. Source scheme follows the one used to build older prototypes, with some changes to increase source efficiency.
To create the electric field the source sends a fast voltage pulse to the electrode. The source was calibrated to know the voltage output in function of parameters set in the instrument, with or without helium gas flow that allows plasma production.
The voltage waveform is a fast positive pulse with a FWHM always lower than $\SI{1}{\micro\second}$. Voltage peak value is linear in a range from \num{4} to \SI{10}{\kilo\volt}, with or without plasma. Pulse repetition rate does not change voltage peak value, it is verified to be the same for rates from \num{5} to \SI{40}{\kilo\hertz}.
With plasma is possible to measure current intensity flowing in a conductive target in front of the source. Also for current is measured a positive pulse with time width lower than $\SI{1}{\micro\second}$. This DBD source is developed for biomedical uses, a fundamental feature is that current deposited on a target should stay below a limit value usually set greater than $\SI{10}{\milli\ampere}$. Measured peak current intensity is always less than \SI{4}{\milli\ampere}. %With measurements done in this thesis it is not possible to exclude a relation between voltage pulse repetition rates and current measured on target. Current peak intensity shows a particular behavior separately for $\num{5} \le f \le \SI{15}{\kilo\hertz}$ and for $\num{20} \le f \le \SI{40}{\kilo\hertz}$, however the number of points is too low to determine the behaviour.
 A useful estimate of the effective current intensity sensed by a living target is an average value in a longer time interval. For a time interval of \SI{1}{\milli\second} the effective current is $\le \SI{0.3}{\milli\ampere}$, below any possible limit value for medical application.
Current peak value results proportional to voltage peak value on the electrode, showing that plasma adds a resistive load between electrode and target. Plasma electrical resistance is estimated between \SI{2.98(11)}{\mega\ohm} and \SI{11.82(219)}{\mega\ohm}, different for different voltage peaks.
In table \ref{tab:el_sum} there is a summary of the electric parameters here described.
\begin{table}[h]
 \centering
 \begin{tabular}{cc}
  \toprule
                    &Range\\
  \midrule
  $f$               &\num{5}-\SI{40}{\kilo\hertz}\\
  $V_{p}$           &\num{4}-\SI{10}{\kilo\volt}\\
  $I_{p}$           &\num{0}-\SI{4}{\milli\ampere}\\
  $I_{\text{eff}}$  &\num{0}-\SI{0.3}{\milli\ampere}\\
  $R_{\text{pl}}$   &\num{3}-\SI{12}{\mega\ohm}\\
  \bottomrule
 \end{tabular}
 \caption{Electric parameters of PCC and produced plasma. $f$ is the voltage pulse repetition rate; $V_{p}$ is the voltage peak applied on the electrode; $I_{p}$ is current peak intensity measured on a conductive target in front of the source; $I_{\text{eff}}$ is the average current intensity in a time interval of \SI{1}{\milli\second}; $R_{\text{pl}}$ is an estimation of plasma resistance.}
 \label{tab:el_sum}
\end{table}


The main analysis done in this work revolvs around the study of plasma propagation. Plasma produced with helium and neon gas propagates with the shape of \emph{bullets}: localized portions of gas that emit light and move with velocities $\ge \SI{10}{\kilo\meter/\second}$.
Bullets produced by PCC are observed with a fast acquisition setup that integrates plasma light emission in an interval of \SI{15}{\nano\second}. Resulting frames are analyzed to study position, dimensions and evaluate velocities of the bullets. Measurements are taken with helium and neon, for different experimental setups, changing voltage peak value, gas flow and target place in front of the source.

Bullet behavior is qualitatively the same for helium and neon gas, with a flow of \SI{2}{\liter/\minute}.
Plasma is observed around the electrode when voltage reaches a value over \SI{5}{\kilo\volt}. After formation, plasma covers all the nozzle area and the bullet moves away from the electrode, towards the end of the nozzle. Once it reaches the end of the nozzle the bullet speeds up and exits rapidly from it, propagating in the air outside. If there is not a target or if there is an insulating target, the bullet propagates in air decelerating until it stops or until the impact with the target.
When an higher voltage pulse is applied on the electrode, bullets travel inside the nozzle with higher velocity and reach higher distances in air. One difference between helium and neon bullets is that the second ones have higher velocities for the same voltage value, but they travel shorter distances in air. Helium plasma bullets, for a voltage peak value between \num{5.5} and \SI{7.5}{\kilo\volt}, reach velocities inside the nozzle from \num{38} to \SI{60}{\kilo\meter/\second}, traveling distances from the end of the nozzle from \num{4} to \SI{23}{\milli\meter}. Neon plasma bullets, for a voltage peak value between \num{4.5} and \SI{6.5}{\kilo\volt}, reach velocities inside the nozzle from \num{35} to \SI{63}{\kilo\meter/\second}, traveling distances from \num{6} to \SI{15}{\milli\meter}.

A model that describes this exact behavior is not known yet. In this work it is excluded that the motion of the bullet is related to propagation of ion waves and it is demonstrated that it is related to electron mobility. A charged particle in an electric field $E$ moves with a drift velocity $v_{d} = \mu E$, where $\mu$ is the particle mobility. The electron mobility, $\mu_e$ varies in function of reactions that happens inside the plasma. Simulation software \emph{Bolsig+} gives an estimation of transport coefficients in a plasma inside an electric field, including $\mu_e$, with settable composition of ionized gas.

The behaviour of electron drift velocity in a gas of helium, neon or a mix of helium and air reflects results found in the experiment. For a gas of helium or neon around $\SI{300}{\kelvin}$, inside an electric field from \num{300} to \SI{700}{\volt/\milli\meter}, at atmospheric pressure, it is found that $\mu_e$ is a near costant value, higher for neon. %compared to Helium. 
Drift velocities $v_{d} = \mu_e E$ are values $\ge \SI{10}{\kilo\meter/\second}$ and decrease linearly for decreasing electric field. This behaviour is comparable to the one found for bullet velocities inside the nozzle, where the electric field decreases moving away from the electrode and with decreasing voltage.
Electron mobility $\mu_e$ is also evaluated for a gas of helium mixed with another gas of $70\%$ nitrogen and $30\%$ oxygen. Different gas composition introduces different reactions inside the gas, changing electron mobility. A gradual increase in the fraction of the second gas allows to simulate the contact of helium with air. For every value of the electric field drift velocities show a peak with low fraction of air in Helium, decreasing once the percentage grows over $20\%$. This behaviour resembles the acceleration of bullets when there is contact between gas flow and air outside the nozzle.
It can be concluded that bullet velocity follows the behaviour of electron drift velocity, proportional to electron mobility. Further simulations or a model for plasma bullets could invesigate this relation to understand how they are related.


For helium bullets it is studied the influence of gas flow in bullet behaviour, going from \num{1} to \SI{4}{\liter/\minute}. Minimum velocity inside the nozzle and distance reached in air are inversely proportional to gas flow, i.e. they are lower with higher flow. Inside the nozzle there is always laminar flow, Reynold numbers are $R_n < \num{200}$, distant from the critical value of \num{2000} identified as the transition value to turbulent flow. Lower velocity, and consequently the shorter distances reached, could be related to the different contact point between helium and air. If the flow is higher, the bullet propagates for more distance in a gas of (almost) pure helium, decelerating more before the acceleration due to impact with air.

The presence of an insulating target in front of the source does not influence bullet propagation. After the impact between bullet and target, both helium and neon bullets continue to propagate on the surface of the insulating target. The propagation on the target is observed as a round shaped figure with increasing diameter on the target. Maximum diameter and the velocity of expansion are proportional to voltage value and inversly proportional to distance between source and target. This behaviour is analogous to bullet propagation in air.

A conductive grounded target changes the behaviour of bullet propagation: velocities in air do not decrease before reaching the target. The bullet stops only after the impact with the target, reaching further distances. A conductive grounded target attracts the bullet. 
With the conductive target it is also possible to compare current measurements with the bullet motion. The most intersting result is that the current pulse starts to rise before the bullet reaches the target. Current pulse's starting time and maximum time are inversly proportional to distance between source and target and are proportional to voltage applied on the electrode. This phenomenon can be explained with the hypotesis that there is a portion of the bullet that does not emit light but carries charge. When an higher voltage is applied on the electrode all the bullet moves with higher velocity and current is measured before.


Bullet formation and propagation are not observed if plasma is produced ionizing argon. From the measurements with the fast camera it is possible to see that argon plasma does not form an homogeneous zone of glowing plasma, a bullet, but inside the nozzle there are several filaments. Those filaments start at the electrode and have different lengths in each frame. During the discharge the average value of the position reached by the end of the filaments moves towards the end of the nozzle. Once this value reachs the end of the nozzle it can be seen that plasma is expelled as formations of round shaped glowing gas, short in height and width.
Those circles have a well defined position and dimensions, allowing to study the motion of the average values of their position and dimensions, as done with position and dimensions of Helium and Neon bullet. The result is that the average value of those formations moves away from nozzle end, with velocities from \num{30} to \SI{60}{\kilo\meter/\second}, depending on the presence of a conductive grounded target or of a grounded ring placed around the nozzle.


Reactive species produced inside plasma are thought to be related to the mechanism behind therapeutic effects of plasma biomedical applications. With spectrometric measurements it is possible to observe which of these species are inside the plasma.
In this study plasma emission is analyzed with two spectrometers in the wavelength range between \num{200} and \SI{800}{\nano\meter}.

Helium plasma emission lines are recognized with an high resolution spectrometer. It is possible to distinguish lines relative to \ce{OH} rotational transitions, \ce{N_2} rotational and vibrational transitions, \ce{N_2^+} vibrational transitions and other lines relative to \ce{O} and \ce{He}.
The ratio between emission lines in a rotational band is related to the rotational temperature of the relative molecule. Rotational temperature can be considered an estimation of gas temperature. Resulting temperatures for \ce{OH} and \ce{N_2} are $T_{r,\ce{OH}} = \SI{352(38)}{\kelvin}$ and $T_{r,\ce{N_2}} = \SI{321(41)}{\kelvin}$, two values compatible with each other and comparable to room temperature. From \ce{N_2} vibrational transitions is possible to give also an estimation of its vibrational temperature $T_{v,\ce{N_2}} = \SI{3405(154)}{\kelvin}$.

Emission is different for different positions inside the plasma plume. With a mini spectrometer is observed emission intensity at different heights of the column of helium plasma, neon plasma and argon plasma. Emission intensity relative to \ce{N_2} always increases going from the electrode to the air outside, while specific emission intensity for the gas used to start the discharge decreases. It is an expected result as the concentration of the gas used to start the discharge is higher inside the nozzle, while concentration of nitrogen is higher in the air outside.

Emission depends also on gas composition. With the same mini spectrometer is observed the variation in emission introducing a percentage of nitrogen or argon in a flux of Helium. Nitrogen up to the $10\%$ of total gas flux does not increase emission intensity for \ce{N_2} lines. Nitrogen is the principal component of air so it is possible that the percentage inserted in the gas flux is low respect the quantity already present.
With an increasing percentage of argon there is an increase in the emission intensity relative to argon lines, and a decrease of the one relative to \ce{N_2} and \ce{He} lines.


A fundamental requirement for plasma produced by PCC is that it cannot increase temperature target over a certain safe limit. In this work is estimated power deposited by plasma on a target. As explained before, plasma production is a pulsed phenomenon, the power relative to a single pulse can be evaluated considering the time interval of interaction between plasma and target of \SI{1}{\micro\second}.
Typical values for the effective power deposited on the target are between \num{20} and \SI{50}{\milli\watt}, while the power relative to a single pulse is between \num{2} and \SI{6}{\watt}.
Power for a single pulse decreases with increasing pulse repetition rate. It is an unexpected result that could imply a different bullet behavior for different pulse repetition rates.
Power for a single pulse increases rapidly with increasing voltage, while it decreases with increasing distance. It is expected because plasma bullets have more speed when there is an higher peak voltage on the electrode or when the target is placed closer to the source.
Under the condition of negligible interaction between heating transfer mechanisms in biological living tissues and plasma deposition, with those power estimations it is possible to evaluate the increase in temperature when plasma produced by PCC will be applied for blood coagulation.


A future development for this thesis is the formulation of a model that explains all the observations here presented. A model that describes plasma bullets would help to understand and predict plasma electric behavior, plasma radiation emission and plasma power deposition on a target. Measurements made in this work give description on DBD plasma phenomenology that will be used to test future models.

\clearpage

\printbibliography

\end{document}
